% Compile this UTF-8 document with XeLaTeX or LuaLaTeX.
\documentclass[11pt]{article}
\usepackage{fontspec}
\usepackage{amsmath,amssymb,mathtools,bm}
\usepackage[margin=1in]{geometry}
\setlength{\parindent}{0pt}
\setlength{\parskip}{0.5em}
\begin{document}
\section*{2 Basic Linear Regression}
{\small Source: Regression Modeling with Actuarial and Financial Applications (Frees)/Regression Modeling with Actuarial and Financial Applications.pdf\par}
\medskip
2
Basic Linear Regression
Chapter Preview. This chapter considers regression in the case of only one explanatory
variable. Despite this seeming simplicity, most of the deep ideas of regression can be
developed in this framework. By limiting ourselves to the one variable case, we are
 able to express many calculations using simple algebra. This will allow us to develop
our intuition about regression techniques by reinforcing it with simple demonstrations.
 Further, we can illustrate the relationships between two variables graphically because we
are working in only two dimensions. Graphical tools prove important for developing a
link between the data and a model.
 2.1 Correlations and Least Squares
Regression is about relationships. Specifically, we will study how two variables,
an \(x\) and a \(y\) , are related. We want to be able to answer questions such as, If
 we change the level of \(x\) , what will happen to the level of \(y\) ? If we compare
two subjects that appear similar except for the \(x\) measurement, how will their \(y\)
 measurements differ? Understanding relationships among variables is critical for
quantitative management, particularly in actuarial science, where uncertainty is
 so prevalent.
It is helpful to work with a specific example to become familiar with key
 concepts. Analysis of lottery sales has not been part of traditional actuarial
 practice, but it is a growth area in which actuaries could contribute.
@ EMPIRICAL
Example: Wisconsin Lottery Sales. State of Wisconsin lottery administrators Filename is
 are interested in assessing factors that affect lottery sales. Sales consists of online "WiscLottery"
lottery tickets that are sold by selected retail establishments in Wisconsin. The
tickets are generally priced at \(\S1 . 0 0\) , so the number of tickets sold equals the
 lottery revenue. We analyze average lottery sales (SALES) over a \(4 0\) -week period,
April 1998 through January 1999, from fifty randomly selected areas identified
by postal (ZIP) code within the state of Wisconsin.
Although many economic and demographic variables might influence sales,
 our first analysis focuses on population (POP) as a key determinant. Chapter 3
will show how to consider additional explanatory variables. Intuitively, it seems
clear that geographic areas with more people have higher sales. So, other things
23
\clearpage
24 Basic Linear Regression
Table 2.1 Summary
Statistics of Each Standard
Variable Variable Mean Median Deviation Minimum Maximum
POP 9,311 4,406 11,098 280 39,098
SALES 6,495 2,426 8,103 189 33,181
Source: Frees and Miller (2003).
Figure 2.1 Frequency Frequency
Histograms of 25 30
 population and sales. 20 25
Each distribution is
skewed to the right. 20
indicating that there 15 15
 are many small areas 10 10
compared to a few
 areas with larger sales 5 5
 and populations. 0 0
0 20,000 40,000 0 15,000 35,000
POP SALES
 being equal, a larger \(x \,=P O P\) means a larger \(y=S A L E S\) . However, the lottery
is an important source of revenue for the state, and we want to be as precise as
 possible.
 A little additional notation will be useful subsequently. In this sample, there are
fifty geographic areas and we use subscripts to identify each area. For example,
\(y_{1}=1 , \! 2 8 5 . 4\) represents sales for the first area in the sample that has population
\(x_{1}=4 3 5\) . Call the ordered pair \(( x_{1} , \; y_{1} ) \,=\, ( 4 3 5 , \; 1 2 8 5 . 4 )\) the first observation.
Extending this notation, the entire sample containing fifty observations may be
represented by \(( x_{1} , \, y_{1} ) , \, \dots, \, ( x_{5 0} , \, y_{5 0} )\) . The ellipses \(( \ldots)\) mean that the pattern is
continued until the final object is encountered. We will often speak of a generic
Begin by working with member of the sample, referring to \(( x_{i} , \, y_{i} )\) as the \(i\) th observation.
each variable Datasets can get complicated, so it will help if you begin by working with
 separately. each variable separately. The two panels in Figure 2.1 show histograms that give
 a quick visual impression of the distribution of each variable in isolation of the
other. Table 2.1 provides corresponding numerical summaries. To illustrate, for
the population variable (POP), we see that the area with the smallest number
contained 280 people, whereas the largest contained 39,098. The average, over
50 Zip codes, was 9,311.04. For our second variable, sales were as low as \$189
and as high as \$33,181.
As Table 2.1 shows, the basic summary statistics give useful ideas of the
structure of key features of the data. After we understand the information in each
variable in isolation of the other, we can begin exploring the relationship between
 the two variables.
\clearpage
 2.1 Correlations and Least Squares 25
SALES 0 Figure 2.2 A scatter
plot of the lottery
 30,000 0 data. Each of the 50
 plotting symbols
25,000 corresponds to a \(\mathrm{Z i p}\)
code in the study. This
 20,000 0 C postal areas with
figure suggests that
15,000 \(\circ\) \(\circ\) P larger populations
have greater lottery
10,000 revenues.
5000 \(\circ8\)
\(\circ\)
0
0 10,000 20,000 30,000 40,000
POP
 Scatter Plot and Correlation Coefficients -- Basic Summary Tools
The basic graphical tool used to investigate the relationship between the two
variables is a scatter plot, such as in Figure 2.2. Although we may lose the
 exact values of the observations when graphing data, we gain a visual impression
of the relationship between population and sales. From Figure 2.2, we see that
 areas with larger populations tend to purchase more lottery tickets. How strong
 is this relationship? Can knowledge of the area's population help us anticipate
the revenue from lottery sales? We explore these two questions here.
One way to summarize the strength of the relationship between two variables
 is through a correlation statistic.
Defnition . The ordinary, or Pearson, correlation coefficient is defined as
\[
\begin{aligned} {r=\frac{1} {( n-1 ) s_{x} s_{y}} \sum_{i=1}^{n} \left( x_{i} \,-\, \overline{{x}} \right) ( y_{i} \,-\, \overline{{y}} ) \, .} \\ \end{aligned}
\]
defined in Section 1.2, with similar notation tor \(s_{x}\) \(s_{y}=\sqrt{( n \,-\, 1 )^{-1} \, \sum_{i=1}^{n} \left( y_{i} \,-\, \overline{{y}} \right)^{2}}\)
Here, we use the sample standard deviation
Although there are other correlation statistics, the correlation coefficient
 devised by Pearson (1895) has several desirable properties. One important prop-
erty is that, for any dataset, \(r\) is bounded by \(- 1\) and 1, that is, \(- 1 \leq r \leq1\) .
(Exercise 2.3 provides steps for you to check this property.) If \(r\) is greater than
zero, the variables are said to be (positively) correlated. If \(r\) is less than zero,
 the variables are said to be negatively correlated. The larger the coefficient is in
 absolute value, the stronger is the relationship. In fact, if \(r=1\) , then the variables
are perfectly correlated. In this case, all of the data lie on a straight line that
 goes through the lower-left- and upper-right-hand quadrants. If \(r=-1 .\) then all
\clearpage
26 Basic Linear Regression
 of the data lie on a line that goes through the upper-left- and lower-right-hand
quadrants. The coefficient \(r\) is a measure of a linear relationship between two
The coeffcient \(r\) is \({}^{a}\) variables.
 measure of a linear The correlation coefficient is said to be location and scale invariant. Thus,
 relationship between ehvarables enerof loca  eaateriea c \(y_{i}\) will increase by \(r\) . For
 two variables. example, if we add \(\S1 0 0\) to the sales of each \(\mathrm{Z i p}\)
100. However, \(\overline{{y}}\) , the average purchase price, will also increase by 100, so that
 the deviation \(y_{i}-\overline{{y}}\) remains unchanged, or invariant. Further, the scale of each
 variable does not matter in the calculation of \(r\) . For example, suppose we divide
each population by 1,000 so that \(x_{i}\) now represents population in thousands.
Thus, \(\overline{{x}}\) is also divided by 1,000 and you should check that \(s_{x}\) is also divided by
1,000. Thus, the standardized version of \(x_{i} , \, ( x_{i} \,-\, \overline{{x}} ) / s_{x}\) , remains unchanged, or
 invariant. Many statistical packages compute a standardized version of a variable
by subtracting the average and dividing by the standard deviation. Now, let's use
\(y_{i , s t d}=( y_{i} \,-\, \overline{{y}} ) / s_{y}\) and \(x_{i , s t d}=( x_{i} \,-\, \overline{{x}} ) / s_{x}\) to be the standardized versions of \(y_{i}\)
and \(x_{i}\) , respectively. With this notation, we can express the correlation coefficient
as \(r \,=\, ( n \,-\, 1 )^{-1} \, \sum_{i=1}^{n} x_{i , s t d} \, \times\, y_{i , s t d} .\)
The correlation coefficient is said to be a dimensionless measure. This is
because we have taken away dollars, and all other units of measures, by consider-
 ing the standardized variables \(x_{i , s t d}\) and \(y_{i , s t d}\) Because the correlation coefficient
does not depend on units of measure, it is a statistic that can readily be compared
 The correlation across different datasets.
coeffcient is location In the world of business, the term correlation is often synonymous with the
and scale imvariant. It  term " relationship? For the purposes of this text, we use the term correlation
is dimensionless.
when referring only to linear relationships. The classic nonlinear relationship
is \(y=x^{2}\) a quadratic relationship. Consider this relationship and the fictitious
 dataset for \(x \, , \, \{-2 , \, 1 , \, 0 , \, 1 , \, 2 \}\) . Now, as an exercise (2.2), produce a rough graph
of the dataset:
\(x_{i}\) \(- 2\) 2 0 1 2
\(i\) 1 3 4 5
\(- 1\)
\(y_{i}\) \(4 ~\) 1 0 1 4
The correlation coefficient for this dataset turns out to be \(r=0\) (check this).
Thus, despite the fact that there is a perfect relationship between \(x\) and \(y ~ (=x^{2} )\) ,
 there is a zero correlation. Recall that location and scale changes are not relevant
in correlation discussions, so we could easily change the values of \(x\) and \(y\) to be
 more representative of a business dataset.
How strong is the relationship between \(y\) and \(x\) for the lottery data? Graph-
ically, the response is a scatter plot, as in Figure 2.2. Numerically, the main
response is the correlation coefficient, which turns out to be \(r \,=0 . 8 8 6\) for this
dataset. We interpret this statistic by saying that SALES and POP are (positively)
correlated. The strength of the relationship is strong because \(r \,=0 . 8 8 6\) is close
to one. In summary, we may describe this relationship by saying that there is a
 strong correlation between SALES and POP.
\clearpage
 2.1 Correlations and Least Squares 27
 Method of Least Squares
Now we begin to explore the question, Can knowledge of population help us
understand sales? To respond to this question, we identify sales as the response
or dependent variable. The population variable, which is used to help understand
 sales, is called the explanatory or independent variable.
 Suppose that we have available the sample data of fifty sales \(\{y_{1} , \, \ldots, \, y_{5 0} \}\)
and your job is to predict the sales of a randomly selected Zip code. Without
 knowledge of the population variable, a sensible predictor is simply \(\overline{{y}}=6 , 4 9 5\)
the average of the available sample. Naturally, you anticipate that areas with
 larger populations will have greater sales. That is, if you also have knowledge of
population, then can this estimate be improved? If so, then by how much?
To answer these questions, the first step assumes an approximate linear rela-
tionship between \(x\) and \(y\) . To fit a line to our data set, we use the method of
 least squares. We need a general technique so that, if different analysts agree on
the data and agree on the fitting technique, then they will agree on the line. If
 different analysts fit a dataset using eyeball approximations, in general, they will
 arrive at different lines, even when using the same dataset.
 The method begins with the line \(y=b_{0}^{*}+b_{1}^{*} x\) , where the intercept and slope,
\(b_{0}^{*}\) and \(b_{1}^{*}\) are merely generic values. For the \(i\) th observation, \(y_{i}-\left( b_{0}^{*}+b_{1}^{*} x_{i} \right)\)
represents the deviation of the observed value \(y_{i}\) from the line at \(x_{i}\) . The
quantity
\[
\SS( b_{0}^{\ast} , \, b_{1}^{\ast} )=\sum_{i=1}^{n} \left( y_{i} \,-\left( b_{0}^{\ast}+b_{1}^{\ast} x_{i} \right) \right)^{2}
\]
 represents the sum of squared deviations for this candidate line. The least squares
 method consists of determining the values of \(b_{0}^{*}\) and \(b_{1}^{*}\) that minimize \(\SS( b_{0}^{\ast} , \, b_{1}^{\ast} )\) .
 This is an easy problem that can be solved by calculus, as follows. Taking partial
derivatives with respect to each argument yields
\[
\frac{\partial} {\partial b_{0}^{*}} S S ( b_{0}^{*} , \, b_{1}^{*} )=\sum_{i=1}^{n} (-2 ) \left( y_{i} \,-\left( b_{0}^{*}+b_{1}^{*} x_{i} \right) \right)
\]
and
\[
\frac{\partial} {\partial b_{1}^{*}} S S ( b_{0}^{*} , b_{1}^{*} )=\sum_{i=1}^{n} (-2 x_{i} ) \left( y_{i}-\left( b_{0}^{*}+b_{1}^{*} x_{i} \right) \right) .
\]
The reader is invited to take second partial derivatives to ensure that we are
minimizing, not maximizing, this function. Setting these quantities equal to zero
 and canceling constant terms yields
\[
\begin{aligned} {\sum_{i=1}^{n} \big( y_{i} \,-\big( b_{0}^{\ast}+b_{1}^{\ast} x_{i} \big) \big)=0} \\ \end{aligned}
\]

\clearpage
28 Basic Linear Regression
and
\[
\sum_{i=1}^{n} x_{i} \left( y_{i}-\left( b_{0}^{*}+b_{1}^{*} x_{i} \right) \right)=0 ,
\]
which are known as the normal equations. Solving these equations yields the
 values of \(b_{0}^{*}\) and \(b_{1}^{*}\) that minimize the sum of squares, as follows.
Defnition. The least squares intercept and slope estimates are
\[
b_{1}=r \frac{s_{y}} {s_{x}} \quad\mathrm{~ a n d ~} \quad b_{0}=\overline{{y}} \,-\, b_{1} \overline{{x}} .
\]
The line that they determine, \(\widehat{y}=b_{0}+b_{1} x\) is called the ftted regression line.
We have dropped the asterisk, or star, notation because \(b_{0}\) and \(b_{1}\) are no longer
 candidate values.
Does this procedure yield a sensible line for our Wisconsin lottery sales?
Earlier, we computed \(r \,=0 . 8 8 6\) From this and the basic summary statis-
tics in Table 2.1, we have \(b_{1}=0 . 8 8 6 ( 8 1 0 3 ) / 1 1 \! {,} 0 9 8=0 . 6 4 7\) and \(b_{0}=6 4 9 5-\)
\(( 0 . 6 4 7 ) 9 3 1 1=4 6 9 . 7 .\) This yields the fitted regression line
\[
\widehat{y}=4 6 9 . 7+( 0 . 6 4 7 ) x \, .
\]
The caret, or "" hat, on top of the \(y\) reminds us that this \(\widehat{y}\) or \(\widehat{S A L E S}\) is a fitted
value. One application of the regression line is to estimate sales for a specific
population, say, \(x \,=\, 1 0 \mathrm{,} 0 0 0\) The estimate is the height of the regression line,
which is \(4 6 9 . 7+( 0 . 6 4 7 ) ( 1 0 , 0 0 0 )=6 9 3 9 . 7 .\)
Example: Summarizing Simulations. Regression analysis is a tool for sum-
marizing complex data. In practical work, actuaries often simulate complicated
financial scenarios; it is often overlooked that regression can be used to summa-
rize relationships of interest.
 To illustrate, Manistre and Hancock (2005) simulated many realizations of
a 10-year European put option and demonstrated the relationship between two
actuarial risk measures, the value-at-risk (VaR) and the conditional tail expec-
 tation (CTE). For one example, these authors examined lognormally distributed
stock returns with an initial stock price of \$100, so that in 10 years the price of
the stock would be distributed as
\[
S ( Z )=1 0 0 \operatorname{e x p} \left( ( . 0 8 ) 1 0+. 1 5 \sqrt{1 0} Z \right) ,
\]
 based on an annual mean return of \(1 0 7_{o}\) , standard deviation of \(1 5 \, \mathcal{7}_{o}\) and the
 outcome from a standard normal random variable \(Z\) . The put option pays the
 difference between the strike price, that will be taken to be \$110 for this example,
and \(S ( Z )\) The present value of this option is
\[
C ( Z )=\mathrm{e}^{-0 . 0 6 ( 1 0 )} \mathrm{m a x} \left( 0 , \, 1 1 0 \,-\, S ( Z ) \right) ,
\]
 based on a \(6 7_{o}\) discount rate.
\clearpage
 2.2 Basic Linear Regression Model 29
E
2 conditional tail
expectation (CTE)
12 versus value at risk
(VaR). Based on
\(n=1 , 0 0 0\)
 from a 10-year
2 European put bond.
Source: Manistre and
2 Hancock (2005).
8
0 \(2\) \(4 ~\) 6 8 10 12
VaR Estimates
To estimate the VaR and CTE, for each \(i\) 1000 i.i.d. standard normal
random variables were simulated and used to calculate 1000 present values,
\(C_{i 1} , \ldots, \, C_{i , 1 0 0 0}\) . The 95th percentile of these present values is the estimate of
the value at risk, denoted as \(V a R_{i}\) : The average of the highest \(5 0 ~ (=( 1-. 0 5 ) \times\)
1,000) of the present values is the estimate of the conditional tail expectation,
denoted as \(C T \, E_{i}\) Manistre and Hancock (2005) performed this calculation
\(i=1 , \dots, 1 , 0 0 0\) times; the result is presented in Figure 2.3. The scatter plot
 shows a strong but not perfect relationship between the \(V a \, R\) and the \(C T E\) , the
 correlation coefficient turns out to be \(r \,=\, 0 . 7 8 2\) .
 2.2 Basic Linear Regression Model
The scatter plot, correlation coefficient, and the fitted regression line are useful
 devices for summarizing the relationship between two variables for a specific
 dataset. To infer general relationships, we need models to represent outcomes of
 broad populations.
 This chapter focuses on a basic linear regression model. The " linearregression"
 part comes from the fact that we fit a line to the data. The " basic"part is because we
 use only one explanatory variable, \(x\) . This model is also known as \(\mathrm{a}\) " simple'linear
regression. This text avoids this language because it gives the false impression
 that regression ideas and interpretations with one explanatory variable are always
 straightforward.
We now introduce two sets of assumptions of the basic model, the observables
and the error representations. They are equivalent, but each will help us as we
later extend regression models beyond the basics:
 Basic Linear Regression Model
Observables Representation Sampling Assumptions
F1.E \(y_{i}=\beta_{0}+\beta_{1} x_{i}\) .
F2. \(\{x_{1} , \dots, \, x_{n} \}\) are nonstochastic variables.
F3. Var \(y_{i}=\sigma^{2}\) .
F4. \(\{y_{i} \}\) are independent random variables.
\clearpage
30 Basic Linear Regression
The observables representation focuses on variables that we can see (or
observe), \(( x_{i} , \, y_{i} )\) Inference about the distribution of \(y\) is conditional on the
observed explanatory variables, so that we may treat \(\{x_{1} , \, \dots, \, x_{n} \}\) as nonstochas-
tic variables (assumption F2). When considering types of sampling mechanisms
for \(( x_{i} , \, y_{i} )\) , it is convenient to think of a stratifed random sampling scheme, where
 values of \(\{x_{1} , \dots, \, x_{n} \}\) are treated as strata, or groups. Under stratified sampling,
for each unique value of \(x_{i}\) , we draw a random sample from a population. To
 illustrate, suppose you are drawing from a database of firms to understand stock
return performance (y) and want to stratify on the basis of the size of the firm.
If the amount of assets is a continuous variable, then we can imagine drawing a
sample of size 1 for each firm. In this way, we hypothesize a distribution of stock
returns conditional on firm asset size.
 As a digression, you will often see reports that summarize results for the " top
50 managers" or the " bestl00 universities," measured by some outcome variable.
In regression applications, make sure that you do not select observations based on
a dependent variable, such as the highest stock return, because this is stratifying
 that is based on the \(y\) , not the \(x\) . Chapter \(6\) will discuss sampling procedures in
 greater detail.
Stratified sampling also provides motivation for assumption \(\mathrm{F 4}\) , the inde-
pendence among responses. One can motivate assumption \(\mathrm{F 1}\) by thinking of
\(( x_{i} , \, y_{i} )\) as a draw from a population, where the mean of the conditional distri-
bution of \(y_{i}\) given \(\{x_{i} \}\) is linear in the explanatory variable. Assumption F3 is
known as homoscedasticity, which we will discuss extensively in Section 5.7. See
Goldberger (1991) for additional background on this representation.
A fifth assumption that is often implicitly used is as follows:
F5. \(\{y_{i} \}\) are normally distributed.
 This assumption is not required for many statistical inference procedures because
central limit theorems provide approximate normality for many statistics of inter-
est. However, formal justification for some, such as \(t\) -statistics, do require this
 additional assumption.
In contrast to the observables representation, an alternative set of assump-
tions focuses on the deviations, or errors, in the regression, defined as \(\varepsilon_{i}=\)
\(y_{i}-( \beta_{0}+\beta_{1} x_{i} )\) :
 Basic Linear Regression Model
Error Representation Sampling Assumptions
El. \(y_{i}=\beta_{0}+\beta_{1} x_{i}+\varepsilon_{i}\)
E2. \(\{x_{1} , \dots, \, x_{n} \}\) are nonstochastic variables.
E3.E \(\varepsilon_{i}=0\) and Var \(\varepsilon_{i}=\sigma^{2}\)
E4. \(\{\varepsilon_{i} \}\) are independent random variables.
 The error representation is based on the Gaussian theory of errors (see Stigler,
1986, for a historical background). Assumption E1 assumes that \(y\) is in part due
\clearpage
2.2 Basic Linear Regression Model 31
Table 2.2 Summary
Data Measures Imneeso iline Variance Sample
 Summary Measures of the
 Population and
 Population Parameters \(\beta_{0}\) \(\beta_{1}\) \(\sigma^{2}\)
Sample Statistics \(b_{0}\) \(b_{1}\) \(s^{2}\)
Figure 2.4 The
distribution of the
True Unknown response varies by the
 Regression Line level of the
 explanatory variable.
The Center of Each Normal
Curve Is at the Height of
 the Regression Line
Each Response Tends
: to Fall Near the Height of
 the Regression Line
\(\mathrm{x_{1}}\) \(\mathrm{x_{2}}\) \(\mathrm{x_{3}}\)
 to a linear function of the observed explanatory variable, \(x\) . Other unobserved
 variables that influence the measurement of \(y\) are interpreted to be included in the
error term, \(\varepsilon_{i}\) , which is also known as the disturbance term. The independence of
errors, \(\mathrm{E 4}\) , can be motivated by assuming that \(\{\varepsilon_{i} \}\) is realized through a simple
 random sample from an unknown population of errors.
 Assumptions E1-- E4 are equivalent to F1--F4. The error representation pro-
 vides a useful springboard for motivating goodness-of-fit measures (Section 2.3).
However, a drawback of the error representation is that it draws the attention
 from the observable quantities \(( x_{i} , \, y_{i} )\) to an unobservable quantity, \(\{\varepsilon_{i} \}\) . To illus-
trate, the sampling basis, viewing \(\{\varepsilon_{i} \}\) as a simple random sample, is not directly
verifiable because one cannot directly observe the sample \(\{\varepsilon_{i} \}\) . Moreover, the
 assumption of additive errors in \(\mathrm{E 1}\) will be troublesome when we consider non-
 linear regression models.
Figure 2.4 illustrates some of the assumptions of the basic linear regression
 model. The data \(( x_{\boldsymbol{1}} , \, y_{\boldsymbol{1}} ) , ( x_{\boldsymbol{2}} , \, y_{\boldsymbol{2}} )\) , and \(( x_{3} , \, y_{3} )\) are observed and are represented by
 the circular opaque plotting symbols. According to the model, these observations
should be close to the regression line \(\textrm{E} y=\beta_{0}+\beta_{1} x\) . Each deviation from the
 line is random. We will often assume that the distribution of deviations may be
represented by a normal curve, as in Figure 2.4.
 The basic linear regression model assumptions describe the underlying pop-
 ulation. Table 2.2 highlights the idea that characteristics of this population can
 be summarized by the parameters \(\beta_{0} , \, \beta_{1}\) , and \(\sigma^{2}\) . In Section 2.1, we summarized
 data from a sample, introducing the statistics \(b_{0}\) and \(b_{1}\) . Section 2.3 will introduce
\(s^{2}\) the statistic corresponding to the parameter \(\sigma^{2}\) .
\clearpage
32 Basic Linear Regression
Figure 2.5
Geometric display of
the deviation \(y\) \(y-\acute{y}\) \(\hat{y}=b_{0}+b_{1} x\)
decomposition.
\(\acute{y}\)
\(\overset{\wedge} {y} \!-\! \overline{{y}}=b_{1} ( x \!-\! \overline{{x}} )\)
\(\overline{{y}}\) \(x-\overline{{x}}\)
\(\overline{{x}}\) \(x\)
 2.3 Is the Model Useful? Some Basic Summary Measures
Although statistics is the science of summarizing data, it is also the art of arguing
with data. This section develops some of the basic tools used to justify the basic
 linear regression model. \(\mathrm{A}\) scatter plot may provide strong visual evidence that
\(x\) influences \(y\) ; developing numerical evidence will enable us to quantify the
strength of the relationship. Further, numerical evidence will be useful when we
consider other datasets where the graphical evidence is not compelling.
2.3.1 Partitioning the Variability
The squared deviations, \(( y_{i} \,-\, \overline{{y}} )^{2}\) , provide a basis for measuring the spread of
the data. If we wish to estimate the \(i\) th dependent variable without knowledge
of \(x\) , then \(\overline{{y}}\) is an appropriate estimate and \(y_{i}-\overline{{y}}\) represents the deviation of
the estimate. We use Total \(S S \,=\, \sum_{i=1}^{n} \left( y_{i} \,-\, \overline{{y}} \right)^{2}\) the total sum of squares, to
represent the variation in all of the responses.
 Suppose now that we also have knowledge of \(x\) , an explanatory variable. Using
 the fitted regression line, for each observation we can compute the corresponding
ftted value, \(\widehat{y}_{i} \,=\, b_{0} \,+\, b_{1} x_{i}\) . The fitted value is our estimate with knowledge of
the explanatory variable. As before, the difference between the response and
the fitted value, \(y_{i} \,-\, \widehat{y}_{i}\) , represents the deviation of this estimate. We now have
two " estimates"of \(y_{i}\) ; these are \(\widehat{y}_{i}\) and \(\overline{{y}}\) . Presumably, if the regression line is
useful, then \(\widehat{y}_{i}\) is a more accurate measure than \(\overline{{y}}\) . To judge this usefulness, we
 algebraically decompose the total deviation as
\[
\begin{array} {c} {\underbrace{y_{i}-\overline{{y}}}_{\mathrm{t o t a l}}} \\ {\mathrm{d e v i a t i o n}} \\ \end{array}=\underbrace{y_{i}-\widehat{y}_{i}}_{\mathrm{d e v i a t i o n}} \begin{array} {c} {+\underbrace{\widehat{y}_{i}-\overline{{y}}}} \\ {+\ \mathrm{e x p l a i n e d}} \\ {\mathrm{d e v i a t i o n}} \\ \end{array} \tag{2.1}
\]
Interpret this equation as " the deviation without knowledge of \(x\) equals the
deviation with knowledge of \(x\) plus the deviation explained by \(x\) " Figure 2.5 is
 a geometric display of this decomposition. In the figure, an observation above
\clearpage
2.3 Is the Model Useful? Some Basic Summary Measures 33
 the line was chosen, yielding a positive deviation from the fitted regression line,
 to make the graph easier to read. \(\mathrm{A}\) good exercise is to draw a rough sketch
 corresponding to Figure 2.5 with an observation below the fitted regression line.
Now, from the algebraic decomposition in equation (2.1), square each side of
 the equation and sum over all observations. After a little algebraic manipulation,
this yields
\[
\begin{aligned} {\sum_{i=1}^{n} \left( y_{i} \,-\, \overline{{y}} \right)^{2}=\sum_{i \,=1}^{n} \left( y_{i} \,-\, \widehat{y_{i}} \right)^{2}+\sum_{i \,=1}^{n} \left( \widehat{y_{i}} \,-\, \overline{{y}} \right)^{2} .} \\ \end{aligned} \tag{2.2}
\]
We rewrite this as Total \(S S=E r r o r ~ S S-\) + Regression \(\mathit{S S}\) , where \(S \, S\) stands for
sum of squares. We interpret
: Total \(\mathit{S S}\) as the total variation without knowledge of \(x\)
: Error \(S \, S\) as the total variation remaining after the introduction of \(x\)
: Regression \(S \, S\) as the difference between Total \(S \, S\) and Error \(\mathit{S S}\) , or the total
variation explained through knowledge of \(x\)
When squaring the right-hand side of equation (2.1), we have the cross-product
term \(2 \, ( y_{i} \,-\, \widehat{y}_{i} \, ) \, ( \widehat{y}_{i} \,-\, \overline{{y}} )\) . With the algebraic manipulation, one can check that the
sum of the cross-products over all observations is zero. This result is not true for
 all fitted lines but is a special property of the least squares fitted line.
In many instances, the variability decomposition is reported through only a
 single statistic.
Defnition. The coeffcient of determination is denoted by the symbol \(R^{2}\) .
called "R-square," and defined as follows:
\[
R^{2}=\cfrac{R e g r e s s i o n \; S S} {T o t a l \; S S} .
\]
We interpret \(R^{2}\) to be the proportion of variability explained by the regression
line. In one extreme case where the regression line fits the data perfectly, we
 have Error \(S S=0\) and \(R^{2}=1 .\) In the other extreme case where the regression
line provides no information about the response, we have Regression \(S S=0\)
and \(R^{2}=0 .\) The coefficient of determination is constrained by the inequalities
\(0 \, \leq\, R^{2} \, \leq\, 1\) I with larger values implying a better fit.
2.3.2 The Size of a Typical Deviation: \(\boldsymbol{s}\)
In the basic linear regression model, the deviation of the response from the regres-
 sion line, \(y_{i}-( \beta_{0}+\beta_{1} x_{i} ) .\) is not an observable quantity because the parameters
\(\beta_{0}\) and \(\beta_{1}\) are not observed. However, by using estimators \(b_{0}\) and \(b_{1}\) , we can
 approximate this deviation using
\[
e_{i} \,=\, y_{i} \,-\, \widehat{y}_{i} \,=\, y_{i} \,-( b_{0} \,+\, b_{1} x_{i} ) \, ,
\]
 known as the residual.
\clearpage
34 Basic Linear Regression
Residuals will be critical to developing strategies for improving model spec-
ification in Section 2.6. We now show how to use the residuals to estimate \(\sigma^{2}\) .
From a first course in statistics, we know that if one could observe the deviations
\(\varepsilon_{i}\) , then a desirable estimate of \(\sigma^{2}\) would be \(( n \,-\, 1 )^{-1} \sum_{i=1}^{n} \, \left( \varepsilon_{i} \,-\, \overline{{\varepsilon}} \right)^{2} .\) Because
\(\{\varepsilon_{i} \}\) are not observed, we use the following.
Defnition . An estimator of \(\sigma^{2}\) , the mean square error (MSE), is defined as
\[
\begin{aligned} {s^{2}=\frac{1} {n-2} \sum_{i=1}^{n} {e_{i}}^{2}} \\ \end{aligned} . \tag{2.3}
\]
The positive square root, \(s=\sqrt{s^{2}}\) is called the residual standard deviation.
Comparing the definitions of \(s^{2}\) and \(( n \,-\, 1 )^{-1} \sum_{i \,=1}^{n} \, ( \varepsilon_{i} \,-\, \overline{{\varepsilon}} )^{2} ,\) you will see two
important differences. First, in defining \(s^{2}\) , we have not subtracted the average
residual from each residual before squaring. This is because the average residual
is zero, a special property of least squares estimation (see Exercise 2.14). This
\(s^{2}\) is an unbiased result can be shown using algebra and is guaranteed for all datasets.
estimator of \(\sigma^{2}\) Second, in defining \(s^{2}\) we have divided by \(n-2\) instead of \(n-1\) . Intuitively,
 dividing by either \(n\) or \(n-1\) tends to underestimate \(\sigma^{2}\) . The reason is that, when
fitting lines to data, we need at least two observations to determine a line. For
example, we must have at least three observations for there to be any variability
about a line. How much " freedom"is there for variability about a line? We will
 say that the error degrees of freedom is the number of observations available, \(n\)
minus the number of observations needed to determine a line, 2 (with symbols,
\(d f=n-2 )\) . However, as we saw in the least squares estimation subsection, we
do not need to identify two actual observations to determine a line. The idea is
 that if an analyst knows the line and \(n-2\) observations, then the remaining two
observations can be determined, without variability. When dividing by \(n-2\) , it
 can be shown that \(s^{2}\) is an unbiased estimator of \(\sigma^{2}\) :
We can also express \(s^{2}\) in terms of the sum of squares quantities. That is,
\[
\begin{aligned} {s^{2}=\frac{1} {n-2} \sum_{i=1}^{n} {( y_{i}-\widehat y_{i} )^{2}}=\frac{E r r o r \ S S} {n-2}=M S E .} \\ \end{aligned}
\]
This leads us to the analysis of variance, or ANOVA, table:
ANOVA Table
Source Sum of Squares \(d f\) Mean Square
Regression Regression \(\mathit{S S}\) \(1 ~\) Regression MS
Error Error \(\mathit{S S}\) \(n-2\) MSE
Total Total \(\mathit{S S}\) \(n-1\)
\clearpage
2.4 Properties of Regression Coeffcient Estimators 35
 The ANOVA table is merely a bookkeeping device used to keep track of the
 sources of variability; it routinely appears in statistical software packages as part
 of the regression output. The mean square column figures are defined to be the
sums of square (SS) figures divided by their respective degrees of freedom \(\left( d f \right)\) .
In particular, the mean square for errors \(( M S E )\) equals \(s^{2}\) , and the regression sum
 of squares equals the regression mean square. This latter property is specific to
 the regression with one variable case; it is not true where we consider more than
 one explanatory variable.
The error degrees of freedom in the ANOVA table is \(n-2\) . The total degrees
of freedom is \(n-1\) , reflecting the fact that the total sum of squares is centered
 about the mean (at least two observations are required for positive variability).
 The single degree of freedom associated with the regression portion means that
the slope, plus one observation, is enough information to determine the line.
This is because it takes two observations to determine a line and at least three
 observations for there to be any variability about the line.
The analysis of variance table for the lottery data is as follows:
ANOVA Table
Source Sum of Squares \(d f\) Mean Square
Regression 2,527,165,015 1 2,527,165,015
Error 690,116,755 48 14,377,432
Total 3,217,281,770 \(4 9\)
From this table, you can check that \(R^{2}=7 8 . 5^{\not\expandafter{\romannumeral1}} \! \! \! \acute{\romannumeral1} \! \! \! \acute{\ G}\) and \(s=3 , 7 9 2 .\)
2.4 Properties of Regression Coefftient Estimators
 The least squares estimates can be expressed as a weighted sum of the responses.
To see this, define the weights
\[
w_{i}=\cfrac{x_{i} \ -\overline{{x}}} {s_{x}^{2} ( n-1 )} .
\]
 Because the sum of \(x\) -deviations : \(( x_{i} \,-\, \overline{{x}} )\) is zero, we see that \(\sum_{i=1}^{n} w_{i}=0\) . Thus,
we can express the slope estimate
\[
\begin{aligned} {b_{1}=r \frac{s_{y}} {s_{x}}=\frac{1} {( n-1 ) s_{x}^{2}} \sum_{i=1}^{n} \left( x_{i}-\overline{{x}} \right) \left( y_{i}-\overline{{y}} \right)=\sum_{i=1}^{n} w_{i} \left( y_{i}-\overline{{y}} \right)=\sum_{i=1}^{n} w_{i} \, y_{i} .} \\ \end{aligned}
\]
4)
 The exercises ask the reader to verify that \(b_{0}\) can also be expressed as a weighted
sum of responses, so our discussion pertains to both regression coefficients.
Because regression coefficients are weighted sums of responses, they can be
 affected dramatically by unusual observations (see Section 2.6).
\clearpage
36 Basic Linear Regression
Regression Because \(b_{1}\) is a weighted sum, it is straightforward to derive the expectation
coeffcients are and variance of this statistic. By the linearity of expectations and Assumption F1,
weighted sums of the we have
responses.
\[
\begin{aligned} {\textrm{E} b_{1}=\sum_{i=1}^{n} w_{i} \textrm{E} y_{i}=\beta_{0} \sum_{i=1}^{n} w_{i}+\beta_{1} \sum_{i=1}^{n} w_{i} x_{i}=\beta_{1} .} \\ \end{aligned}
\]
Tht is. \(b_{1}\) is an unbiased estimator of \(\beta_{1}\) Here, the sum \(\sum_{i=1}^{n} \, w_{i} \, x_{i} \ \ =\\)
\(\left[ s_{x}^{2} ( n \,-\, 1 ) \right]^{-1} \sum_{i \,=1}^{n} \left( x_{i} \,-\, \overline{{x}} \right) x_{i} \ =\left[ s_{x}^{2} ( n \,-\, 1 ) \right]^{-1} \sum_{i \,=1}^{n} \left( x_{i} \,-\, \overline{{x}} \right)^{2}=1 .\) From the
definition of the weights, some easy algebra also shows that \(\sum_{i=1}^{n} w_{i}^{2}=\)
\(1 / \left( s_{x}^{2} ( n-1 ) \right)\) . Further, the independence of the responses implies that the vari-
ance of the sum is the sum of the variances, and thus we have
\[
\begin{aligned} {\mathrm{V a r ~} b_{1}=\sum_{i=1}^{n} w_{i}^{2} \mathrm{V a r ~} y_{i}=\frac{\sigma^{2}} {s_{x}^{2} ( n-1 )} .} \\ \end{aligned}
\]
Replacing \(\sigma^{2}\) by its estimator \(s^{2}\) and taking square roots leads to the following.
 Defnition . The standard error of \(b_{1}\) the estimated standard deviation of \(b_{1}\) .
 is defined as
\[
s e ( b_{1} )=\frac{s} {s_{x} \sqrt{n \mathrm{~}-\mathrm{~} 1}} . \tag{2.5}
\]
This is our measure of the reliability, or precision, of the slope estimator. Using
equation (2.5), we see that \(s e ( b_{1} )\) is determined by three quantities, \(n , \, s\) , and \(s_{x}\) :
A standard error is an as follows:
estimated standard
 deviation. - If we have more observations so that \(n\) becomes larger, then \(s e ( b_{1} )\) becomes
 smaller, other things equal.
: If the observations have a greater tendency to lie closer to the line so that \(s\)
becomes smaller, then \(s e ( b_{1} )\) becomes smaller, other things equal.
: If values of the explanatory variable become more spread out so that \(s_{x}\)
increases, then \(s e ( b_{1} )\) becomes smaller, other things equal.
 Smaller values of \(s e ( b_{1} )\) offer a better opportunity to detect relations between
\(y\) and \(x\) . Figure 2.6 illustrates these relationships. Here, the scatter plot in the
 middle has the smallest value of \(s \, e ( b_{1} )\) Compared with the middle plot, the left-
 hand plot has a larger value of \(s\) and thus \(s e ( b_{1} )\) Compared with the right-hand
 plot, the middle plot has a larger \(s_{x}\) and thus smaller value of \(s e ( b_{1} )\)
Equation (2.4) also implies that the regression coefficient \(b_{1}\) is normally dis-
tributed. That is, recall from mathematical statistics that linear combinations of
normal random variables are also normal. Thus, if Assumption \(\mathrm{F 5}\) holds, then
\(b_{1}\) is normally distributed. Moreover, several versions of central limit theorems
\clearpage
2.5 Statistical Inference 37
Figure 2.6 These
three scatter plots
\(\circ\) \(\circ_{\mathrm{O}}\) \(\circ\) \(\circ\) 08 00 C \(x\) . The
Iinearelaionshir
\(y\)
\(\circ\) \(\circ^{\circ}\) \(0_{0_{0}}\) \% plot on the left
C exhibits greater
variability about the
line than the plot in
the middle. The plot
\(x\) \(x\) \(x\) on the right exhibits a
smaller standard
 deviation in \(x\) than the
exists for weighted sums (see,e. Serfling, 1980).Thus, as discussed in Sectionplotinthe midle.
1.4, if the responses \(y_{i}\) are even approximately normally distributed, then it will
 be reasonable to use a normal approximation for the sampling distribution of
\(b_{1}\) . Using \(s e ( b_{1} )\) as the estimated standard deviation of \(b_{1}\) , for large values of
\(n\) we have that \(\left( b_{1}-\beta_{1} \right) / s e ( b_{1} )\) has an approximate standard normal distribu-
tion. Although we will not prove it here, under Assumption F5 \(\left( b_{1}-\beta_{1} \right) / s e ( b_{1} )\)
follows a \(t\) -distribution with degrees of freedom \(d f=n-2\) .
 2.5 Statistical Inference
Having fit a model with a dataset, we can make a number of important statements.
Generally, it is useful to think about these statements in three categories: (i) tests
of hypothesized ideas, (ii) estimates of model parameters, and (ii) predictions of
new outcomes.
2.5.1 Is the Explanatory Variable Important? The t-Test
We respond to the question of whether the explanatory variable is important
by investigating whether \(\beta_{1}=0\) The logic is that if \(\beta_{1}=0\) then the basic
linear regression model no longer includes an explanatory variable \(x\) . Thus,
 we translate our question of the importance of the explanatory variable into a
 narrower question that can be answered using the hypothesis-testing framework.
This narrower question is, IS \(H_{0} : \beta_{1}=0\) valid? We respond to this question by
looking at the test statistic:
\[
\mathrm{t-r a t i o}=\frac{\mathrm{e s t i m a t o r-h y p o t h e s i z e d ~ v a l u e ~ o f ~ p a r a m e t e r}} {\mathrm{s t a n d a r d ~ e r r o r ~ o f ~ t h e ~ e s t i m a t o r}} .
\]
provides additional
Appendix A3.3
details about the
t-distribution,
 the hypothesized value of \(\beta_{1}\) is \(\; \; \; \; \; \; \; \; \; \; \; \; \; \; \; \; \; \; \; \; \; \; \; \; \; \; \; \; \; \; \; \; \; \; \; \; \; \; \; \; \; \; \; \; \; \; \; \; \; \; \; \; \; \; \; \; \; \; \; \; \; \; \; \; \; \; \; \; \; \; \; \; \; \; \; \; \; \; \; \; \; \; \; \; \; \; \; \; \; \; \; \; \; \; \; \; \; \; \; \; \; \; \; \; \; \; \; \; \; \; \; \; \; \; \; \; \; \; \; \; \; \; \; \; \; \; \; \; \; \; \; \; \; \; \; \; \; \; \; \; \; \; \; \; \; \; 0 \; \; \; \; \; 0\) , we examine the \(t\) -ratio \(t ( b_{1} )=b_{1} / s e ( b_{1} )\) because including a graph and
For the case of \(H_{0} : \beta_{1}=0\) . This is the appropriate standardization because,distribution table.
under the null hypothesis and the model assumptions described in Section \(2 . 4\)
the sampling distribution of \(t ( b_{1} )\) can be shown to be the \(t\) -distribution with \(d f=\)
\(n-2\) - Thus, to test the null hypothesis \(H_{0}\) against the alternative \(H_{a} : \beta_{1} \neq0\)
\clearpage
38 Basic Linear Regression
Table 2.3
Decision-Making Alternative Hypothesis \(( H_{a} )\) Procedure: Reject \(H_{0}\) in favor of \(H_{a}\) if
Procedures for
Testing \(H_{0} : \beta_{1}=d\) \(\beta_{1} > d\)
\[
\begin{aligned} {t \mathrm{-r a t i o}} & {{} > t_{n-2 , 1-\alpha} .} \\ {t \mathrm{-r a t i o}} & {{} <-t_{n-2 , 1-\alpha} .} \\ {| t \mathrm{-r a t i o} |} & {{} > t_{n-2 , 1-\alpha/ 2} .} \\ \end{aligned}
\]
\(\beta_{1} < d\)
\(\beta_{1} \neq d\)
Note: The significance level is \(\alpha\) . Here, \(t_{n-2 , 1-\alpha}\) is the \(( 1-\alpha)\) th percentile
 from the \(t\) -distribution using \(d f=n \,-2\) degrees of freedom. The test statistic
is \(t\) -ratio \(= ( b_{1}-d ) / s e ( b_{1} )\) .
we reject \(H_{0}\) in favor of \(H_{a}\) if \(\left| t ( b_{1} ) \right|\) exceeds a \(t\) -value. Here, this \(t\) -value is a
 percentile from the \(t\) -distribution using \(d f=n-2\) . We denote the significance
level as \(\alpha\) and this \(t\) -value as \(t_{n-2 , 1-\alpha/ 2}\)
Example: Lottery Sales, Continued. For the lottery sales example, the residual
 standard deviation is \(s \,=3 7 9 2\) From Table 2.1, we have \(s_{x}=1 1 , 0 9 8\) . Thus, the
standard error of the slope is \(s e ( b_{1} )=3 7 9 2 / ( {1 1} {\mathrm{,} 0 9 8} \sqrt{5 0 \,-\, 1} )=0 . 0 4 8 8\) . From
Section 2.1, the slope estimate is \(b_{1}=0 . 6 4 7\) . Thus, the \(t\) -statistic is \(t ( b_{1} )=\)
\(0 . 6 4 7 / 0 . 0 4 8 8=1 3 . 4\) We interpret this by saying that the slope is 13.4 standard
errors above zero. For the significance level, we use the customary value of
\(\alpha\,=\, 5^{o} \! /_{o}\) . The 97.5th percentile from a \(t\) -distribution with \(d f=5 0-2=4 8\)
 degrees of freedom is \(t_{4 8 , 0 . 9 7 5}=2 . 0 1 1\) . Because \(| 1 3 . 4 | > 2 . 0 1 1\) , we reject the
 null hypothesis that the slope \(\beta_{1}=0\) in favor of the alternative that \(\beta_{1} \neq0\)
 Making decisions by comparing a \(t\) -ratio to a \(t\) -value is called a \(t \mathrm{-} t e s t\) . Testing
\(H_{0} : \beta_{1}=0\) versus \(H_{a} : \beta_{1} \neq0\) is just one of many hypothesis tests that can
be performed, although it is the most common. Table 2.3 outlines alternative
 decision-making procedures. These procedures are for testing \(H_{0} : \beta_{1}=d\) , where
\(d\) is a user-prescribed value that may be equal to zero or any other known value.
For example, in our Section 2.7 example, we will use \({d=1}\) to test financial
 theories about the stock market.
Alternatively, one can construct probability \(\left( p \right)\) values and compare these to
given significant levels. The \(p\) -value is a useful summary statistic for the data
 analyst to report because it allows the report reader to understand the strength of
 the deviation from the null hypothesis. Table 2.4 summarizes the procedure for
 calculating \(p\) -values.
 Another interesting way of addressing the question of the importance of an
explanatory variable is through the correlation coefficient. Remember that the
 correlation coefficient is a measure of linear relationship between \(x\) and \(y\) . Let's
denote this statistic by \(r ( y , \, x )\) This quantity is unaffected by scale changes
in either variable. For example, if we multiply the \(x\) variable by the number
\(b_{1}\) , then the correlation coefficient remains unchanged. Further, correlations are
\clearpage
2.5 Statistical Inference 39
Table 2.4
Alternative Probability Values for
Hypothesis ( \(H_{a}\) \(\beta_{1} > d\) \(\beta_{1} < d\) \(\beta_{1} \neq d\) Testing \(H_{0} : \beta_{1}=d\)
\(p\) -Value \(\operatorname* {P r} ( t_{n-2} \, > \, t \mathrm{-r a t i o} )\) \(\operatorname* {P r} ( t_{n-2} < t \mathrm{-r a t i o} )\) \(\operatorname* {P r} ( | t_{n-2} | \, > \, | t\) -ratio()
Notes: Here, \(t_{n-2}\) is a \(t\) -distributed random variable with \(d f=n-2\) degrees of free-
dom. The test statistic is \(\quad t \mathrm{-r a t i o}=( b_{1}-d ) / s e ( b_{1} )\) .
unchanged by additive shifts. Thus, if we add a number, say, \(b_{0}\) , to each \(x\)
 variable, then the correlation coefficient remains unchanged. Using a scale change
and an additive shift on the \(x\) variable can be used to produce the fitted value
\(\widehat{y}=b_{0}+b_{1} x\) . Thus, using notation, we have \(r ( y , \, x )=\vert r ( y , \, \widehat{y} ) \vert\) . We may thus
 interpret the correlation between the responses and the explanatory variable to be
 equal to the correlation between the responses and the fitted values. This leads,
 then, to the following interesting algebraic fact, \(R^{2}=r^{2}\) . That is, the coefficient
 of determination equals the correlation coefficient squared. This is much easier to
 interpret if one thinks of \(r\) as the correlation between observed and fitted values.
See Exercise 2.13 for steps useful in confirming this result. \(R^{2}=r^{2} .\)
2.5.2 Confidence Intervals
 Investigators often cite the formal hypothesis-testing mechanism to respond to the
 question, Does the explanatory variable have a real influence on the response? \(\mathrm{A}\)
 natural follow-up question is, To what extent does \(x\) affect \(y\) ? To a certain degree,
 one could respond using the size of the \(t\) -ratio or the \(p\) -value. However, in many
instances, a confdence interval for the slope is more useful.
To introduce confidence intervals for the slope, recall that \(b_{1}\) is our point
 estimator of the true, unknown slope \(\beta_{1}\) . Section \(2 . 4\) argued that this estimator has
 standard error \(s e ( b_{1} )\) and that \(\left( b_{1}-\beta_{1} \right) / s e ( b_{1} )\) follows a \(t\) -distribution with \(n-2\)
degrees of freedom. Probability statements can be inverted to yield confidence
intervals. Using this logic, we have the following confidence interval for the
 slope \(\beta_{1}\) :
Defnition . A \(1 0 0 ( 1-\alpha)^{\sigma_{o}}\) confidence interval for the slope \(\beta_{1}\) is
\[
b_{1} \pm t_{n-2 , 1-\alpha/ 2} ~ s e ( b_{1} ) . \tag{2.6}
\]
As with hypothesis testing, \(t_{n-2 , 1-\alpha/ 2}\) is the \(( 1-\alpha/ 2 ) \mathrm{t h}\) percentile from the
\(t\) -distribution with \(d f=n \,-2\) degrees of freedom. Because of the two-sided
nature of confidence intervals, the percentile is \(1-( 1\) - confidence level)/2. In
this text, for notational simplicity, we generally use a 95\% confidence interval,
 so the percentile is \(1-( 1-. 0 . 9 5 ) / 2=0 . 9 7 5\) . The confidence interval provides
 a range of reliability that measures the usefulness of the estimate.
\clearpage
40 Basic Linear Regression
In Section 2.1, we established that the least squares slope estimate for the
lottery sales example is \(b_{1}=0 . 6 4 7\) The interpretation is that if a Zip code's
 population differs by 1,000, then we expect mean lottery sales to differ by \$647.
How reliable is this estimate? It turns out that \(s e ( b_{1} )=0 . 0 4 8 8\) and thus an
 approximate 95\% confidence interval for the slope is
\[
0 . 6 4 7 \pm( 2 . 0 1 1 ) ( . 0 4 8 8 ) ,
\]
or (0.549, 0.745). Similarly, if population differs by 1,000, a 95\% confidence
interval for the expected change in sales is (549, 745). Here, we use the \(t\) -value
\(t_{4 8 , 0 . 9 7 5}=2 . 0 1 1\) because there are \(4 8 ~ (=n ~-2 )\) degrees of freedom and, for a
\(9 5 7_{o}\) confidence interval, we need the 97.5th percentile.
2.5.3 Prediction Intervals
In Section 2.1, we showed how to use least squares estimators to predict the lottery
sales for a Zip code, outside of our sample, with a population of 10,000. Because
prediction is such an important task for actuaries, we formalize the procedure so
 that it can be used on a regular basis.
 To predict an additional observation, we assume that the level of explanatory
variable is known and is denoted by \(x_{*}\) . For example, in our previous lottery sales
example, we used \(x_{*}=1 0 \mathrm{,} 0 0 0\) . We also assume that the additional observation
 follows the same linear regression model as the observations in the sample.
Using our least square estimators, our point prediction is \(\widehat{y}_{\ast}=b_{0}+b_{1} x_{\ast}\) , the
 height of the fitted regression line at \(x_{*}\) . We may decompose the prediction error
into two parts:
pre
\[
\underbrace{y_{*}-\widehat{y}_{*}}_{\mathrm{d i c t i o n ~ e r r o r}} \,=\, \underbrace{\beta_{0}-b_{0}+( \beta_{1}-b_{1} ) \, x_{*}}_{\mathrm{r e g r e s s i o n ~ l i n e ~ a t ~} x_{*}} \,+\, \underbrace{\varepsilon_{*}}_{\mathrm{r e v i a t i o n ~ o f ~ t h e ~ a d}} .
\]
Iditional
s mean
It can be shown that the standard error of the prediction is
\[
s e ( p r e d )=s \sqrt{1+\frac{1} {n}+\frac{\left( x_{*}-\overline{{x}} \right)^{2}} {\left( n-1 \right) s_{x}^{2}}} .
\]
As with \(s e ( b_{1} )\) the terms \(n^{-1}\) and \(\left( x_{*}-\overline{{x}} \right)^{2} / \left[ ( n-1 ) s_{x}^{2} \right]\) become close to zero
 as the sample size \(n\) becomes large. Thus, for large \(n\) , we have that \(s e ( p r e d ) \approx s\)
reflecting that the error in estimating the regression line at a point becomes
 negligible and deviation of the additional response from its mean becomes the
entire source of uncertainty.
\clearpage
2.6 Building a Better Model: Residual Analysis 41
 Defnition . A \(1 0 0 ( 1-\alpha)^{\sigma_{o}}\) prediction interval at \(x_{*}\) is
\[
\widehat{y}_{*} \pm\ t_{n-2 , 1-\alpha/ 2} \ s e ( p r e d ) , \tag{2.7}
\]
 where the \(t\) -value \(t_{n-2 , 1-\alpha/ 2}\) is the same as used for hypothesis testing and the
 confidence interval.
For example, the point prediction at \(x_{\ast}=1 0 {,} 0 0 0 \, \mathrm{i s} \, \, \widehat{y}_{\ast}=4 6 9 . 7+0 . 6 4 7 \, ( 1 0 {,} 0 0 0 )=\)
6939.7. The standard error of this prediction is
\[
\begin{aligned} {s e ( p r e d )} & {{}=3 7 9 2 \sqrt{1+\frac{1} {5 0}+\frac{\left( 1 0 {,} 0 0 0-9 {,} 3 1 1 \right)^{2}} {( 5 0-1 ) ( 1 1 {,} 0 9 8 )^{2}}}} \\ \end{aligned}=3 8 2 9 . 6 .
\]
With a \(t\) -value equal to 2.011, this yields an approximate 95 \(\mathcal{J}_{o}\) prediction interval:
\[
6 9 3 9 . 7 \pm( 2 . 0 1 1 ) ( 3 8 2 9 . 6 )=6 9 3 9 . 7 \pm7 , 7 0 1 . 3=(-7 6 1 . 6 , \, \, \, 1 4 , 6 4 1 . 0 ) .
\]
We interpret these results by first pointing out that our best estimate of lottery
 sales for a \(\mathrm{Z i p}\) code with a population of 10,000 is \$6,939.70. Our 95\% prediction
 interval represents a range of reliability for this prediction. If we could see many
Zip codes, each with a population of 10,000, on average we would expect that
about \(1 9\) out of 20, or \(9 5 \surd_{o}\) , would have lottery sales between \(\S0\) and \$14,641.
It is customary to truncate the lower bound of the prediction interval to zero if
 negative values of the response are deemed to be inappropriate.
2.6 Building a Better Model: Residual Analysis
Quantitative disciplines calibrate models with data. Statistics takes this one step
further, using discrepancies between the assumptions and the data to improve
model specification. We will examine the Section 2.2 modeling assumptions in
light of the data and use any mismatch to specify a better model; this process
is known as diagnostic checking (like when you go to a doctor and he or she
 performs diagnostic routines to check your health). Statistics uses
We will begin with the Section \(2 . 2\) error representation. Under this set ofdiscrepancies between
assumptions, the deviations \(\{\varepsilon_{i} \}\) are identically and independently distributed the assumptions and
the data to improve
(i.i.d) and, under assumption F5, normally distributed. To assess the validity of model specifcation.
these assumptions, one uses (observed) residuals \(\{e_{i} \}\) as approximations for the
(unobserved) deviations \(\{\varepsilon_{i} \}\) . The basic theme is that if the residuals are related to
 a variable or display any other recognizable pattern, then we should be able to take
 advantage of this information and improve our model specification. The residuals If the residuals are
 should contain little or no information and represent only natural variation from related to a variable
 the sampling that cannot be attributed to any specific source. Residual analysis or display any other.
 recognizable pattern,
is the exercise of checking the residuals for patterns. then we should be
If detected, the discrepancies can be corrected with the appropriate adjustments of this information
 and improve our
 model specifcation.
\clearpage
42 Basic Linear Regression
 Model Misspecifcation Issues
(i) Lack of Independence. There may exist relationships among the devia-
tions \(\{\varepsilon_{i} \}\) so that they are not independent.
(ii) Heteroscedasticity. Assumption E3 that indicates that all observations
have a common (though unknown) variability, known as homoscedas-
ticity. Heteroscedascity is the term used when the variability varies by
observation.
(iii) Relationships between Model Deviations and Explanatory Variables.
 If an explanatory variable has the ability to help explain the deviation \(\varepsilon\) ,
 then one should be able to use this information to better predict \(y\) .
(iv) Nonnormal Distributions. If the distribution of the deviation represents
 a serious departure from normality, then the usual inference procedures
 are no longer valid.
(v) Unusual Points. Individual observations may have a large effect on the
regression model fit, meaning that the results may be sensitive to the
 impact of a single observation.
This list will serve you throughout your study of regression analysis. Of course,
 with only an introduction to basic models, we have not yet seen alternative models
that might be used when we encounter such model discrepancies. In this book's
Part \(\amalg\) on time series models, we will study lack of independence among data
 ordered over time. Chapter 5 will consider heteroscedasticity in further detail. The
 introduction to multiple linear regression in Chapter 3 will be our first look at han-
 dling relationships between \(\{\varepsilon_{i} \}\) and additional explanatory variables. We have,
however, already had an introduction to the effect of normal distributions, seeing
that \(q \, q\) plots can detect nonnormality and that transformations can help induce
 approximate normality. In this section, we discuss the effects of unusual points.
 Much of residual analysis is done by examining a standardized residual, a
residual divided by its standard error. An approximate standard error of the resid-
 ual is \(s\) ; in Chapter 3, we will give a precise mathematical definition. There are
two reasons we often examine standardized residuals in lieu of basic residu-
als. First, if responses are normally distributed, then standardized residuals are
 approximately realizations from a standard normal distribution. This provides a
reference distribution to compare values of standardized residuals. For example,
if a standardized residual exceeds two in absolute value, this is considered unusu-
ally large and the observation is called an outlier. Second, because standardized
residuals are dimensionless, we get carryover of experience from one dataset to
 another. This is true regardless of whether the normal reference distribution is
 applicable.
 Outliers and High Leverage Points. Another important part of residual analy-
sis is the identification of unusual observations in a dataset. Because regression
estimates are weighted averages with weights that vary by observation, some
\clearpage
2.6 Building a Better Model: Residual Analysis 43
 Variables 19 Base Points \(\mathrm{A}\) \(\mathbf{B}\) CTabi 2.5 19 a
Points Plus Three
\(x\) 1.51.72.02.22.52.52.72.93.03.5 \(3 . 4\) \(9 . 5\) 9.5Typesof Umusual
Observations
\(y\) 3.02.53.53.03.13.63.23.94.04.08.08.02.5
\(x\) 3.84.24.34.64.05.15.15.25.5
\(y\) 4.24.14.8 4.25.15.15.14.85.3
9 Figure 2.7 Scatter
8 oA oB plot of 19 base plus 3 \(\mathrm{C}\) .
7 unusual points,
labeled A, B and
6
y 5
\(2 ~\) : \(\circ\underset{\circ} {\underbrace{\circ}} \overset{\otimes} {\otimes}\) 0000 oC
\(4 ~\)
3 0
0 \(2 ~\) \(4 ~\) 6 8 10
\(x\)
observations are more important than others. This weighting is more important
 than many users of regression analysis realize. In fact, the example here demon-
strates that a single observation can have a dramatic effect in a large dataset.
There are two directions in which a data point can be unusual, the horizontal
and vertical directions. \(\mathrm{B y}\) - unusual, I mean that an observation under consid-
eration seems to be far from the majority of the dataset. An observation that is
unusual in the vertical direction is called an outlier. An observation that is unusual
 in the horizontal directional is called a high leverage point. An observation may
 be both an outlier and a high leverage point.
@ EMPIRICAL
ExpePo
\(1 9\) points plus three points, labeled \(\mathrm{A , \ B}\) , and \(\mathrm{C}\) given in Figure 2.7 and"OutlierExample"
 Table 2.5. Think of the first nineteen points as " good"observations that represent
 some type of phenomena. We want to investigate the effect of adding a single
aberrant point.
To investigate the effect of each type of aberrant point, Table 2.6 summarizes
the results of four separate regressions. The first regression is for the nineteen
base points. The other three regressions use the nineteen base points plus each
 type of unusual observation.
Table 2.6 shows that a regression line provides a good fit for the nineteen
base points. The coefficient of determination, \(R^{2}\) , indicates that about 89\% of
 the variability has been explained by the line. The size of the typical error, \(s\) , is
about 0.29, small compared to the scatter in the \(y\) -values. Further, the \(t\) -ratio for
 the slope coefficient is large.
\clearpage
44 Basic Linear Regression
Table 2.6 Results.
 from Four Data \(b_{0}\) \(b_{1}\) \(s\) \(R^{2} ( \mathcal{V}_{\! o} )\) \(t ( b_{1} )\)
Regressions
19 Base Points 1.869 0.611 0.288 89.0 11.71
19 Base Points \(+ \mathrm{A}\) 1.750 0.693 0.846 53.7 4.57
19 Base Points \(+ \mathrm{B}\) 1.775 \(0 . 6 4 0\) 0.285 94.7 18.01
19 Base Points \(+ \mathrm{C}\) 3.356 0.155 0.865 10.3 1.44
When the outlier point \(\mathrm{A}\) is added to the nineteen base points, the situation
 deteriorates dramatically. The \(R^{2}\) drops from 89\% to 53.7\%, and \(s\) increases from
about 0.29 to about 0.85. The fitted regression line itself does not change that
much even though our confidence in the estimates has decreased.
An outlier is unusual in the \(y\) -value, but "- unusualin the \(y\) -value" depends on
 the \(x\) -value. To see this, keep the \(y\) -value of point \(\mathrm{A}\) the same but increase the
\(x\) -value and call the point \(\mathrm{B}\) .
When the point \(\mathrm{B}\) is added to the nineteen base points, the regression line
provides a better fit. Point \(\mathrm{B}\) is close to being on the line of the regression fit
 generated by the nineteen base points. Thus, the fitted regression line and the size
of the typical error, \(s\) , do not change much. However, \(R^{2}\) increases from 89\%
 to nearly 95 percent. If we think of \(R^{2}\) as \(1-( E r r o r ~ S S ) / ( T o t a l ~ S S ) ,\) by adding
point \(\mathrm{B}\) we have increased Total \(S \, S\) , the total squared deviations in the \(y^{\prime}\) s even
 though leaving Error SS relatively unchanged. Point \(\mathrm{B}\) is not an outlier, but it is
a high leverage point.
To show how influential this point is, drop the \(y\) -value considerably and call
this the new point \(\mathrm{C}\) When this point is added to the nineteen base points,
the situation deteriorates dramatically. The \(R^{2}\) coefficient drops from 89\% to
\(1 0 7_{o}\) , and the \(s\) more than triples, from 0.29 to 0.87. Further, the regression line
coefficients change dramatically.
Most users of regression at first do not believe that one point in twenty can
have such a dramatic effect on the regression fit. The fit of a regression line can
 always be improved by removing an outlier. If the point is a high leverage point 
and not an outlier, it is not clear whether the fit will improve when the point is
removed.
Simply because you can dramatically improve a regression fit by omitting an
 observation does not mean you should always do so! The goal of data analysis is
to understand the information in the data. Throughout the text, we will encounter
 many datasets where the unusual points provide some of the most interesting
information about the data. The goal of this subsection is to recognize the effects
 of unusual points; Chapter \(5\) will provide options for handling unusual points in
 your analysis.
All quantitative disciplines, such as accounting, economics, linear program-
\(\operatorname* {m i n g}\) , and so on, practice the art of sensitivity analysis. Sensitivity analysis is a
 description of the global changes in a system due to a small local change in an
\clearpage
2.6 Building a Better Model: Residual Analysis 45
Table 2.7
Data \(b_{0}\) \(b_{1}\) \(s\) \(R^{2} ( \ \% )\) \(t ( b_{1} )\) Regression Results
with and without
 With Kenosha \(4 6 9 . 7\) 0.662 3792 \(7 8 . 5\) \(1 3 . 2 6\) Kenosha
\(0 . 6 4 7\)
 Without Kenosha \(- 4 3 . 5\) 2728 88.3 18.82
SALES Figure 2.8 Scatter
30,000 0 plot of SALES versus
0 POP, with the outlier
corresponding to
 25,000 Kenosha Kenosha marked.
20,000 0 C \(\circ\)
15,000 \(\circ\) 0 \(\circ\) 0
10,000 \(\circ\) 0
5000 \(\circ8\)
0
0 10000 20000 30000 40000
POP
element of the system. Examining the effects of individual observations on the
regression fit is a type of sensitivity analysis.
Example: Lottery Sales, Continued. Figure 2.8 exhibits an outlier; the point in
 the upper-left-hand side of the plot represents a \(\mathrm{Z i p}\) code that includes Kenosha,
Wisconsin. Sales for this Zip code are unusually high given its population.
Kenosha is close to the Illinois border; residents from Illinois probably par-
ticipate in the Wisconsin lottery, thus effectively increasing the potential pool of
 sales in Kenosha. Table 2.7 summarizes the regression fit both with and without
this Zip code.
For the purposes of inference about the slope, the presence of Kenosha does
not alter the results dramatically. Both slope estimates are qualitatively similar
 and the corresponding \(t\) -statistics are very high, well above cutoffs for statistical
 significance. However, there are dramatic differences when assessing the quality
 of the fit. The coefficient of determination, \(R^{2}\) increased from 78.5\% to 88.3\%
when deleting Kenosha. Moreover, our typical deviation \(s\) dropped by more
than \$1,000. This is particularly important if we want to tighten our prediction
intervals.
To check the accuracy of our assumptions, it is also customary to check the
normality assumption. One way of doing this is with the \(q \, q\) plot, introduced
\clearpage
46 Basic Linear Regression
Figure 2.9 The \(q q\) Sample Quantiles Sample Quantiles
plots of Wisconsin 0 8,000 o
lottery residuals. The 15,000 0 2,000 \(\infty\)
left-hand panel is 6,000 0
 based on all 50 points. 10,000 4,000
The right-hand panel
 is based on 49 points,
residuals from a 5,000 \(\infty^{\circ}\) 0
regression after
removing Kenosha. 0 2,000
5,000 o 0 \(\omicron\circ\) 6,000 \(\circ\) 0 \(\omicron^{\mathrm{c}}\)
4,000
0
T T
2 1 0 1 2 2 1 0 1 2
Theoretical Quantiles Theoretical Quantiles
in Section 1.2. The two panels in Figures 2.9 are \(q \, q\) plots with and without
 the Kenosha Zip code. Recall that points close to linear indicate approximate
normality. In the right-hand panel of Figure 2.9, the sequence does appear to be
 linear, so residuals are approximately normally distributed. This is not the case in
 the left-hand panel, where the sequence of points appears to climb dramatically for
large quantiles. The interesting thing is that the nonnormality of the distribution
 is due to a single outlier, not to a pattern of skewness that is common to all the
observations.
2.7 Application: Capital Asset Pricing Model
In this section, we study a financial application, the capital asset pricing model
(CAPM). The name is something of a misnomer in that the model is really about
returns based on capital assets, not on the prices themselves. The types of assets
that we examine are equity securities that are traded on an active market, such
as the New York Stock Exchange (NYSE). For a stock on the exchange, we can
relate returns to prices through the following expression:
Return = Price at the end of a period + dividends -- price at the begining of a period
 price at the beginning of a period
If we can estimate the returns that a stock generates, then knowledge of the
 price at the beginning of a generic financial period allows us to estimate the value
at the end of the period (ending price plus dividends). Thus, we follow standard
practice and model returns of a security.
 An intuitively appealing idea, and one of the basic characteristics of the CAPM,
 is that there should be a relationship between the performance of a security and the
market. One rationale is simply that if economic forces are such that the market
improves, then those same forces should act on an individual stock, suggesting
 that it also improve. As noted previously, we measure performance of a security
\clearpage
2.7 Application: Capital Asset Pricing Model 47
 through the return. To measure performance of the market, several market indices
exist that summarize the performance of each exchange. We will use the equally
weighted index of the Standard \(\&\) Poor's 500. The S\&P500 is the collection of
the 500 largest companies traded on the NYSE, where " lage" is identified by
Standard \(\&\) Poor's, a financial services rating organization. The equally weighted
 index is defined by assuming that a portfolio is created by investing \$1.00 in each
 of the 500 companies.
 Another rationale for a relationship between security and market returns comes
from financial economics theory. This is the CAPM theory, attributed to Sharpe
(1964) and Lintner (1965) and based on the portfolio diversification ideas of
 Markowitz (1959). Other things equal, investors would like to select a return with
 a high expected value and a low standard deviation, the latter being a measure of
risk. One of the desirable properties of using standard deviations as a measure
of riskiness is that it is straightforward to calculate the standard deviation of a
portfolio. One needs to know only the standard deviation of each security and
the correlations among securities. A notable security is a risk-free one, that is,
a security that theoretically has a zero standard deviation. Investors often use a
 30-day U.S. Treasury bill as an approximation of a risk-free security, arguing that
 the probability of default of the U.S. government within 30 days is negligible.
Positing the existence of a risk-free asset and some other mild conditions, under
the CAPM theory there exists an efficient frontier called the securities market
\(l i n e\) This frontier specifies the minimum expected return that investors should
 demand for a specified level of risk. To estimate this line, we can use the equation
\[
\textrm{E} r=\beta_{0}+\beta_{1} r_{m} ,
\]
where \(r\) is the security return and \(r_{m}\) is the market return. We interpret \(\beta_{1} r_{m}\) as a
 measure of the amount of security return attributed to the behavior of the market.
Testing economic theory, or models arising from any discipline, involves
collecting data. The CAPM theory is about ex ante (before-the-fact) returns even
though we can only test with ex post (after-the-fact) returns. Before the fact,
the returns are unknown and there is an entire distribution of returns. After the
fact, there is only a single realization of the security and market return. Because
 at least two observations are required to determine a line, CAPM models are
estimated using security and market data gathered over time. In this way, several
 observations can be made. For the purposes of our discussions, we follow standard
 practice in the securities industry and examine monthly prices.
 Data @ EMPIRICAL
To illustrate, consider monthly returns over the five-year period from January Filename is
1986 to December 1990, inclusive. Specifically, we use the security returns from"CAPM"
the Lincoln National Insurance Corporation as the dependent variable (y) and
the market returns from the index of the S\&P500 as the explanatory variable
\(( x )\) At the time, the Lincoln was a large, multiline insurance company, with
headquarters in the Midwest, specifically in Fort Wayne, Indiana. Because it was
\clearpage
48 Basic Linear Regression
Table 2.8 Summary
Statistics of 60 Standard
Monthly Observations Mean Median Deviation Minimum Maximum
LINCOLN 0.0051 0.0075 0.0859 0.2803 0.3147
MARKET 0.0074 0.0142 0.0525 0.2205 0.1275
Source: Center for Research on Security Prices, University of Chicago.
Figure 2.10 Time Monthly Return
 series plot of returns 0.3 LINCOLN]
from Lincoln National 0.2 MARKET
Corporation and the
market. There are 60
monthly returns over 0.1
 the period January
193-Decembe 0.0
0.1
0.2
0.3
1986 1987 1988 1989 1990 1991
Year
well known for its prudent management and stability, it is a good company to
 begin our analysis of the relationship between the market and an individual stock.
We begin by interpreting some of basic summary statistics, in Table 2.8, in
terms of financial theory. First, an investor in Lincoln will be concerned that
the five-year average return, \(\overline{{y}}=0 . 0 0 5 1 0\) is less than the return of the mar-
ket, \(\overline{{x}}=0 . 0 0 7 4 1\) . Students of interest theory recognize that monthly returns can
be converted to an annual basis using geometric compounding. For example,
 the annual return of Lincoln is \(( 1 . 0 0 5 1 )^{1 2}-1=0 . 0 6 2 9 4 6\) or roughly 6.29\%.
This is compared to an annual return of \(9 . 2 6 \not{0} \ (=( 1 0 0 ( ( 1 . 0 0 7 4 1 )^{1 2}-1 ) )\) for the
market. A measure of risk, or volatility, that is used in finance is the standard devi-
ation. Thus, interpret \(\boldsymbol{s}_{y} \,=\, 0 . 0 8 5 9 > \, 0 . 0 5 2 5 4=\boldsymbol{s}_{x}\) to mean that an investment in
Lincoln is riskier than the whole market. Another interesting aspect of Table 2.8
 is that the smallest market return, \(- 0 . 2 2 0 5 2\) is 4.338 standard deviations below
 its average ) \(\left( (-0 . 2 2 0 5 2 \,-\, 0 . 0 0 7 4 1 ) / 0 . 0 5 2 5 4=-4 . 3 3 8 \right)\) This is highly unusual
with respect to a normal distribution.
We next examine the data over time, as is given graphically in Figure 2.10,
which shows scatter plots of the returns versus time, called time series plots. In
Figure 2.10, one can clearly see the smallest market return, and a quick glance at
 the horizontal axis reveals that this unusual point is in October 1987, the time of
 the well-known market crash.
The scatter plot in Figure 2.11 graphically summarizes the relationship
 between Lincoln's return and the return of the market. The market crash is clearly
\clearpage
2.7 Application: Capital Asset Pricing Model 49
LINCOLN Figure 2.11 Scatter
0.4 plot of Lincoln's
return versus the S\&P
0.3 500 return. The
0.2 regression line is
 superimposed,
0.1 ( enabling us to identify
the market crash and
0.0 C 1990 outliers two outliers.
October 1987 crash
0.1 \(\circ\) 00
0.2
0.3 \(\circ\)
0.3 0.2 0.1 0.0 0.1 0.2
MARKET
evident in Figure 2.11 and represents a high leverage point. With the regression
line (described subsequently) superimposed, the two outlying points that can be
seen in Figure 2.10 are also evident. Despite these anomalies, the plot in Fig-
ure 2.11 does suggest that there is a linear relationship between Lincoln and
 market returns.
 Unusual Points
To summarize the relationship between the market and Lincoln's return, a regres-
sion model was fit. The fitted regression is
\[
\widehat{L I N C O L N}=-0 . 0 0 2 1 4+0 . 9 7 3 M A R K E T .
\]
 The resulting estimated standard error, \(s \,=\, 0 . 0 6 9 6\) is lower than the standard
 deviation of Lincoln's returns, \(s_{y} \,=\, 0 . 0 8 5 9\) Thus, the regression model explains
some of the variability of Lincoln's returns. Further, the \(t\) -statistic associated
with the slope \(b_{1}\) turns out to be \(t ( b_{1} )=5 . 6 4 .\) which is significantly large. One
 disappointing aspect is that the statistic \(R^{2}=3 5 . 4^{\mathcal{G}_{0}}\) can be interpreted to mean
 that the market explains only slightly more than a third of the variability. Thus,
even though the market is clearly an important determinant, as evidenced by the
high \(t\) -statistic, it provides only a partial explanation of the performance of the
Lincoln's returns.
In the context of the market model, we may interpret the standard deviation
of the market, \(s_{x}\) , as nondiversifable risk. Thus, the risk of a security can be
decomposed into two components, the diversifiable component and the market
component, which is nondiversifiable. The idea is that, by combining several
securities, we can create a portfolio of securities that, in most instances, will
reduce the riskiness of our holdings when compared with a single security.
Again, the rationale for holding a security is that we are compensated through
higher expected returns by holding a security with higher riskiness. To quantify
\clearpage
50 Basic Linear Regression
 the relative riskiness, it is not hard to show that
\[
s_{y}^{2}=b_{1}^{2} s_{x}^{2}+s^{2} \cfrac{n-2} {n-1} . \tag{2.8}
\]
The riskiness of a security results from the riskiness due to the market plus
the riskiness due to a diversifiable component. Note that the riskiness due to the
 market component, \(s_{x}^{2}\) , is larger for securities with larger slopes. For this reason,
investors think of securities with slopes \(b_{1}\) greater than one as " aggressive" and
slopes less than one as " defensive."
Sensitivity Analysis
The foregoing summary immediately raises two additional issues. First, what
is the effect of the October 1987 crash on the fitted regression equation? We
 know that unusual observations, such as the crash, may potentially influence the
fit a great deal. To this end, the regression was rerun without the observation
 corresponding to the crash. The motivation for this is that the October 1987 crash
represents a combination of highly unusual events (the interaction of several
 automated trading programs operated by the large stock brokerage houses) that
we do not want to represent using the same model as our other observations.
Deleting this observation, the fitted regression is
\[
\widehat{L I N C O L N}=-0 . 0 0 1 8 1 \,+\, 0 . 9 5 6 M A R K E T ,
\]
with \(R^{2}=2 6 . 4^{\mathcal{V}_{0}}\) . \(t ( b_{1} )=4 . 5 2 , \, \, \, s=0 . 0 7 0 2\) ,and \(s_{y} \,=\, 0 . 0 8 1 1\) . We interpret
these statistics in the same fashion as the fitted model including the October
1987 crash. It is interesting to note, however, that the proportion of variability
explained actually decreases when excluding the influential point. This serves
to illustrate an important point. High leverage points are often looked on with
 dread by data analysts because they are, by definition, unlike other observations
 in the dataset and require special attention. However, when fitting relationships
 among variables, they also represent an opportunity because they allow the data
analyst to observe the relationship between variables over broader ranges than
otherwise possible. The downside is that the relationships may be nonlinear or
follow an entirely different pattern when compared to the relationships observed
 in the main portion of the data.
The second question raised by the regression analysis is what can be said
about the unusual circumstances that gave rise to the unusual behavior of
Lincoln's returns in October and November 1990. A useful feature of regres-
sion analysis is to identify and raise the question; it does not resolve it. Because
the analysis clearly pinpoints two highly unusual points, it suggests that the data
 analyst should go back and ask some specific questions about the sources of the
data. In this case, the answer is straightforward. In October 1990, the Travelers'
Insurance Company, a competitor, announced that it would take a large write-
 off in its real estate portfolio because of an unprecedented number of mortgage
defaults. The market reacted quickly to this news, and investors assumed that
\clearpage
2.8 Illustrative Regression Computer Output 51
other large publicly traded life insurers would also soon announce large write-
 offs. Anticipating this news, investors tried to sell their portfolios of, for example,
Lincoln's stock, thus causing the price to plummet. However, it turned out that
investors overreacted to the news and that Lincoln's portfolio of real estate was
indeed sound. Thus, prices quickly returned to their historical levels.
2.8 Illustrative Regression Computer Output
Computers and statistical software packages that perform specialized calculations
 play a vital role in modern-day statistical analyses. Inexpensive computing capa-
 bilities have allowed data analysts to focus on relationships of interest. Specifying
 models that are attractive merely for their computational simplicity is much less
 important now than in times before the widespread availability of inexpensive
computing. An important theme of this text is to focus on relationships of interest
 and to rely on widely available statistical software to estimate the models that we
specify.
With any computer package, generally the most difficult parts of operating
the package are the input, using the commands, and interpreting the output. You
will find that most modern statistical software packages accept spreadsheet or
text-based files, which makes input of data relatively easy. Personal computer
 statistical software packages have menu-driven command languages with easily
accessible on-line help facilities. Once you decide what to do, finding the right
commands is relatively easy.
This section provides guidance in interpreting the output of statistical pack-
ages. Most statistical packages generate similar output. Below, three examples
of standard statistical software packages, EXCEL, SAS, and \(\mathrm{R}\) are given. The
 annotation symbol " [.]"'marks a statistical quantity that is described in the legend.
 Thus, this section provides a link between the notation used in the text and output
 from some of the standard statistical software packages.
EXCEL Output
Regression Statistics
Multiple \(\mathtt{R}\) 0.886283 [d]
\(\mathtt{R}\) Square 0.785497 [k]
Adjusted R Square 0.781028 [1]
Standard Error 3791.758 [j]
Observations 50[a]
 ANOVA
df Ss MS \(\mathtt{F}\) Significance \(\mathtt{F}\)
Regression 1[m] 2527165015 [p] 2527165015 [s] 175.773[u] 1.15757E-17[v]
Residual 48[n] 690116754.8[q] 14377432.39[t]
Total 49[o] 3217281770 [r]
Coefficients Standard Error t stat P-value
Intercept 469.7036[b] 702.9061896[d] 0.668230846[f] 0.507187[h]
\(\mathtt{X}\) Variable 1 0.647095[c] 0.048808085[e] 13.25794257[g] 1.16E-17[i]
\clearpage
52 Basic Linear Regression
SAS Output
The SAS System
The REG Procedure
Dependent Variable: SALEs
Analysis of Variance
Sum of Mean
Source DF Squares Square F Value Pr > F
Model 1[m] 2527165015[p] 2527165015 [s] 175.77[u]<.0001[v]
Error 48[n] 6901l6755[q] 14377432 [t]
Corrected Total 49[o] 3217281770 [r]
Root MSE 3791.75848 [j] R-Square 0.7855[k]
Dependent Mean 6494.82900[e] Adj R-Sq 0.7810[1]
Coeff Var 58.38119 [F]
Parameter Estimates
Parameter Standard
Variable Label DF Estimate Error t Value Pr>It|
Intercept Intercept 1 469.70360[b] 702.90619[d] 0.67[f] 0.5072 [h]
POP POP 1 0.64709[c] 0.04881[e] 13.26 [g] <.0001[i]
R Output
Analysis of Variance Table
Response: SALES
Df Sum Sq Mean Sq F value Pr(>F)
POP 1[m] 2527165015[p] 2527165015[s] 175.77304 [u] <2.22e-16 [v]***
Residuals 48[n] 6901l6755[q] 14377432 [t]
Call: lm(formula = SALES \textasciitilde{} POP)
Residuals:
Min 1Q Median 3Q Max
-6047 -1461 -670 486 18229
Coefficients:
Estimate Std. Error t value Pr \(( > \mid{\mathtt t} \mid\) )
(Intercept) 469.7036[b] 702.9062[c] 0.67[f] 0.51 [h]
POP 0.6471[c] 0.0488 [e] 13.26[g] <2e-16 ***[i]
Signif. codes: 0 *** 0.001 ** 0.01 * 0.05 . 0.1 1
Residual standard error: 3790[j] on 48[n] degrees of freedom
Multiple R-Squared: 0.785[k], Adjusted R-squared: 0.781[l]
\(\mathtt{F}\) -statistic: 176[u] on 1[m] and 48[n] DF, p-value: <2e-16[v]
\clearpage
2.8 Illustrative Regression Computer Output 53
Legend Annotation Defnition, Symbol
[a] Number of observations \(n\) .
[b] The estimated intercept \(b_{0}\)
[c] The estimated slope \(b_{1}\) .
[d] The standard error of the intercept, \(s e ( b_{0} )\)
[e] The standard error of the intercept, \(s \, e ( b_{1} )\)
[f] The \(t\) -ratio associated with the intercept, \(t ( b_{0} )=b_{0} / s e ( b_{0} )\) .
[g] The \(t\) -ratio associated with the slope, \(t ( b_{1} )=b_{1} / s e ( b_{1} ) .\)
[h] The \(p\) -value associated with the intercept; here, \(p \,-\, v a l u e=\, P \, r ( | t_{n-2} | >\)
\(\left| t ( b_{0} ) \right| )\) , where \(t ( b_{0} )\) is the realized value (0.67 here) and \(t_{n-2}\) has a
\(t\) -distribution with \(d f=n \,-2\) .
[i] The \(p\) -value associated with the slope; here, \(p \,-\, v a l u e=\, P \, r ( | t_{n-2} | >\)
\(\left| t ( b_{1} ) \right| )\) , where \(t ( b_{1} )\) is the realized value (13.26 here) and \(t_{n-2}\) has a
\(t\) -distribution with \(d f=n \,-2\) .
[j] The residual standard deviation, \(s\) .
[k] The coefficient of determination, \(R^{2}\)
[l] The coefficient of determination adjusted for degrees of freedom, \(R_{a}^{2}\) (This
term will be defined in Chapter 3.)
[m] Degree of freedom for the regression component: 1, for one explanatory
variable.
[n] Degree of freedom for the error component, \(n-2 ,\) for regression with one
 explanatory variable.
[o] Total degrees of freedoms, \(n-1 .\)
[p] The regression sum of squares, Regression \(S \, S\)
[q] The error sum of squares, Error \(S \, S\) .
[r] The total sum of squares, Total Ss.
[s] The regression mean square, Regression M S = Regression SS/1, for one
 explanatory variable.
[t] The error mean square, \(s^{2}=E r r o r\) MS = Error \(S S / ( n-2 )\) , for one
 explanatory variable.
[u] The F-ratio = (Regression MS)/(Error MS). (This term will be defined
 in Chapter 3.)
[v] The \(p\) -value associated with the \(\boldsymbol{F}\) -ratio. (This term will be defined in
Chapter 3.)
[w] The observation number, \(i\) :
[x] The value of the explanatory variable for the \(i\) th observation, \(x_{i}\)
[y] The response for the \(i\) th observation, \(y_{i}\) .
[z] The fitted value for the \(i\) th observation, \(\widehat{y_{i}}\) .
[A] The standard error of the fit, \(s e ( \widehat{y}_{i} )\) .
[B] The residual for the \(i\) th observation, \(e_{i}\) .
[C] The standardized residual for the \(i\) th observation, \(e_{i} / s \, e ( e_{i} )\) The standard
error \(s \, e ( e_{i} )\) will be defined in Section 5.3.1.
[D] The multiple correlation coefficient is the square root of the coefficient of
 determination, \(R=\sqrt{R^{2}}\) . This will be defined in Chapter 3.
\clearpage
54 Basic Linear Regression
[E] The average response, \(\overline{{y}}\)
[F] The coefficient of variation of the response is \(s_{y} / \overline{{y}} .\) SAS prints out \(1 0 0 s_{y} / \overline{{y}}\)
2.9 Further Reading and References
Relatively few applications of regression are basic in the sense that they use
only one explanatory variable; the purpose of regression analysis is to reduce
 complex relationships among many variables. Section 2.7 described an important
exception to this general rule, the CAPM finance model; see Panjer et al. (1998)
for additional actuarial descriptions of this model. Campbell et al. (1997) give a
financial econometrics perspective.
 Chapter References
Anscombe, Frank (1973). Graphs in statistical analysis. American Statistician 27, 17-- 21.
Campbell, John \(\mathrm{Y .}\) , Andrew W. Lo, and A. Craig MacKinlay (1997). The Econometrics of
Financial Markets. Princeton University Press, Princeton, New Jersey.
 Frees, Edward W., and Tom W. Miller (2003). Sales forecasting using longitudinal data models.
International Journal of Forecasting 20, 97- 111.
Goldberger, Arthur (1991). A Course in Econometrics. Harvard University Press, Cambridge,
Massachusetts.
 Koch, Gary J. (1985). A basic demonstration of the \([-1 , 1 ]\) range for the correlation coefficient.
American Statistician 39, 201- 2.
Lintner, J. (1965). The valuation of risky assets and the selection of risky investments in stock
 portfolios and capital budgets. Review of Economics and Statistics 47, no. 1, 13-- 37.
Manistre, B. John, and Geoffrey H. Hancock (2005). Variance of the CTE estimator. North
American Actuarial Journal 9, no. 2, 129- 56.
Markowitz, Harry (1959). Portfolio Selection: Effcient Diversifcation of Investments. John
Wiley and Sons, New York.
Panjer, Harry H., Phelim P. Boyle, Samuel H. Cox, Daniel Dufresne, Hans U. Gerber, Heinz H.
 Mueller, Hal W. Pedersen, Stanley R. Pliska, Michael Sherris, Elias S. Shiu, and Ken S. Tan
(1998). Financial Economics: With Applications to Investment, Insurance and Pensions.
Society of Actuaries, Schaumburg, Illinois.
Pearson, Karl (1895). Royal Society Proceedings 58, 241.
 Serfling, Robert J. (1980). Approximation Theorems of Mathematical Statistics. John Wiley
and Sons, New York.
Sharpe, William F. (1964). Capital asset prices: A theory of market equilibrium under condi-
tions of risk. Journal of Finance 19, no. 3, 425-- 42.
Stigler, Steven M. (1986). The History of Statistics: The Measurement of Uncertainty before
1900. Harvard University Press, Cambridge, Massachusetts.
2.10 Exercises
Sections 2.1-- 2.2
 2.1. Consider the following dataset
\(i\) 1 2 3
\(x_{i}\) 2 \(- 6\) 7.
\(y_{i}\) 3 4 6
\clearpage
2.10 Exercises 55
Fit a regression line using the method of least squares. Determine \(r , \, b_{1}\) .
and \(b_{0}\) .
2.2. A Perfect Relationship, yet Zero Correlation. Consider the quadratic
relationship \(y=x^{2}\) , with data
\(i\) 1 234 5
\(x_{i}\) \(- 2\) \(- 1\) 0 12.
\(y_{i}\) \(4\) 10 1 4
 a. Produce a rough graph for this dataset.
b. Check that the correlation coefficient is \(r=0\) .
2.3. Boundedness of the Correlation Coeffcient. Use the following steps to
 show that \(r\) is bounded by \(- 1\) and 1 (These steps are due to Koch, 1990).
a. Let \(a\) and \(c\) be generic constants. Verify
\[
\begin{aligned} {0} & {{} \leq\frac{1} {n-1} \sum_{i=1}^{n} \left( a \frac{x_{i}-\overline{{x}}} {s_{x}}-c \frac{y_{i}-\overline{{y}}} {s_{y}} \right)^{2}} \\ {} & {{}=a^{2}+c^{2}-2 a c r .} \\ \end{aligned}
\]
b. Use the results in part (a) to show \(2 a c ( r \,-\, 1 ) \leq( a \,-\, c )^{2}\) .
c. By taking \(a=c ,\) use the result in part (b) to show \(r \, \leq\, 1\) .
d. By taking \(a=-c\) , use the results in part (b) to show \(r \geq-1\) .
e. Under what conditions is \(r=-1\) ? Under what conditions is \(r=1 ?\)
2.4. Regression Coeffcients Are Weighted Sums. Show that the intercept
term, \(b_{0}\) , can be expressed as a weighted sum of the dependent variables.
That is, show that \(b_{0}=\sum_{i=1}^{n} w_{i , 0} y_{i}\) . Further, express the weights in terms
of the slope weights, \(w_{i}\) .
2.5. Another Expression for the Slope as a Weighted Sum.
 a. Using algebra, establish an alternative expression
\[
b_{1}={\frac{\sum_{i=1}^{n} {w e i g h t_{i} °} \times s l o p e_{i}} {\sum_{i=1}^{n} {w e i g h t_{i}}}} .
\]
Here, slope; is the slope between \(( x_{i} , \, y_{i} )\) and \(( \bar{x} , \, \bar{y} )\) . Give a precise form
 for the weight weight: as a function of the explanatory variable \(x\) .
b. Suppose that \(\bar{x} \,=4 , \, \bar{y} \,=3 , \, x_{1}=2\) and \(y_{1}=6 .\) Determine the slope and
weight for the first observation, that is, \(s l o p e_{1}\) and weight .
2.6. Consider two variables, \(y\) and \(x\) . Do a regression of \(y\) on \(x\) to get a slope
coefficient that we call \(b_{1 , x , y}\) . Do another regression of \(x\) on \(y\) to get a
slope coefficient that we call \(b_{1 , \, y , \, x}\) . Show that the correlation coefficient
between \(x\) and \(y\) is the geometric mean of the two slope coefficients up to
sign; that is, show that \(\left\vert r \right\vert=\sqrt{b_{1 , \, x \, , \, y} b_{1 , \, y , \, x}}\) .
2.7. Regression through the Origin. Consider the model \(y_{i}=\beta_{1} x_{i} \,+\, \varepsilon_{i}\) , that
 is, regression with one explanatory variable without the intercept term. This
model is called regression through the origin because the true regression
line \(\textrm{E} y=\beta_{1} x\) passes through the origin (the point \(( 0 , 0 )\) . For this model,
\clearpage
56 Basic Linear Regression
the least squares estimate of \(\beta_{1}\) is that number \(b_{1}\) that minimizes the sum
 of squares \(\mathrm{S S} ( b_{1}^{*} )=\sum_{i=1}^{n} \big( y_{i} \,-\, b_{1}^{*} x_{i} \big)^{2}\) .
a. Verify that
\[
b_{1}=\frac{\sum_{i=1}^{n} x_{i} \, y_{i}} {\sum_{i=1}^{n} x_{i}^{2}} .
\]
b. Consider the model \(y_{i}=\beta_{1} z_{i}^{2}+\varepsilon_{i}\) a quadratic model passing through
the origin. Use the result from part (a) to determine the least squares
estimate of \(\beta_{1}\) .
2.8. a. Show that
\[
\begin{aligned} {s_{y}^{2}=\frac{1} {n-1} \sum_{i=1}^{n} \left( y_{i}-\overline{{y}} \right)^{2}=\frac{1} {n-1} \left( \sum_{i=1}^{n} y_{i}^{2}-n \, \overline{{y}}^{2} \right) .} \\ \end{aligned}
\]
b. Follow the same steps to show \(\sum_{i \mathop{=} 1}^{n} \left( y_{i} \,-\, \overline{{y}} \right) ( x_{i} \,-\, \overline{{x}} )=\sum_{i \mathop{=} 1}^{n} x_{i} \, y_{i} \,-\)
\(n \overline{{x}} ~ \overline{{y}}\) .
c. Show that
\[
b_{1}=\frac{{\sum_{i=1}^{n} \left( y_{i} \,-\, \overline{{y}} \right) \left( x_{i} \,-\, \overline{{x}} \right)}} {{\sum_{i=1}^{n} \left( x_{i} \,-\, \overline{{x}} \right)^{2}}} .
\]
 d. Establish the commonly used formula
\[
b_{1}=\frac{\sum_{i=1}^{n} x_{i} \, y_{i} \,-n \overline{{x}} \, \, \overline{{y}}} {\sum_{i=1}^{n} x_{i}^{2} \,-n \overline{{x}}^{2}} .
\]
2.9. Interpretation of Coeffcients Associated with a Binary Explanatory
Variable. Suppose that \(x_{i}\) only takes on the values \(\; \; \; \; \; \; \; \; \; \; \; \; \; \; \; \; \; \; \; \; \; \; \; \; \; \; \; \; \; \; \; \; \; \; \; \; \; \; \; \; \; \; \; \; \; \; \; \; \; \; \; \; \; \; \; \; \; \; \; \; \; \; \; \; \; \; \; \; \; \; \; \; \; \; \; \; \; \; \; \; \; \; \; \; \; \; \; \; \; \; \; \; \; \; \; \; \; \; \; \; \; \; \; \; \; \; \; \; \; \; \; \; \; \; \; \; \; \; \; \; \; \; \; \; \; \; \; \; \; \; \; \; \; \; \; \; \; \; \; \; \; \; \; \; \; \; \; 0 \; \; \; \; 0\) and 1. Out of the
\(n\) observations, \(n_{1}\) take on the value \(x=0\) The \(n_{1}\) observations have an
 average \(y\) value of \(\overline{{y}}_{1}\) The remaining \(n-n_{1}\) observations have value \(x=1\)
@ EMPIRICAL and an average \(y\) value of \(\overline{{y}}_{2}\) . Use Exercise 2.8 to show that \(b_{1}=\overline{{y}}_{2}-\overline{{y}}_{1}\) .
Filename is 2.10. Nursing Home Utilization. This exercise considers nursing home data
"WiscNursingHome" provided by the Wisconsin Department of Health and Family Services
(DHFS) and described in Exercise 1.2.
 Part 1: Use cost-report year 2000 data, and do the following analysis.
a. Correlations
\(\mathrm{a ( i )}\) .Calculate the correlation between TPY and LOGTPY. Comment
on your result.
a(ii). Calculate the correlation among TPY, NUMBED, and
SQRFOOT. Do these variables appear highly correlated?
 a(iii). Calculate the correlation between TPY and NUMBED/10. Com-
ment on your result.
b.Scatter plots. Plot TPY versus NUMBED and TPY versus SQRFOOT.
Comment on the plots.
c. Basic linear regression.
\(\mathrm{c ( i )}\) . Fit a basic linear regression model using TPY as the outcome
variable and NUMBED as the explanatory variable. Summarize
the fit by quoting the coefficient of determination, \(R^{2}\) , and the
\(t\) -statistic for NUMBED.
\clearpage
2.10 Exercises 57
c(ii). Repeat \(\mathrm{c ( i )}\) using SQRFOOT instead of NUMBED. In terms of
\(R^{2}\) , which model fits better?
c(iii). Repeat \(\mathrm{c ( i )}\) , using LOGTPY for the outcome variable and
LOG(NUMBED) as the explanatory variable.
\(\mathrm{c ( i v )}\) . Repeat c(iii) using LOGTPY for the outcome variable and
LOG(SQRFOOT) as the explanatory variable.
Part 2: Fit the model in Part \(\mathrm{1 . c} ( 1 )\) using 2001 data. Are the patterns
 stable over time?
Sections 2.3-- 2.4
2.11. Suppose that, for a sample size of \(n=3\) you have \(e_{2}=2 4\) and \(e_{3}=-1\) .
Determine \(e_{1}\) .
2.12. Suppose that \(r \,=0 , \, n \,=\, 1 5\) , and \(s_{y}=1 0\) . Determine \(\mathit{S}\) .
2.13. The Correlation Coeffiient and the Coeffiient of Determination. Use
the following steps to establish a relationship between the coefficient of
 determination and the correlation coefficient.
a. Show that \(\widehat{y}_{i} \,-\, \overline{{y}}=b_{1} ( x_{i} \,-\, \overline{{x}} )\)
c. Use part (b) to establish \(R^{2}=r^{2}\) \(S S=\sum_{i=1}^{n} \left( \widehat{y}_{i} \,-\, \overline{{y}} \right)^{2}=b_{1}^{2} s_{x}^{2} ( n \,-\, 1 ) .\)
b. Use part (a) to show that Regress
2.14. Show that the average residual is zero, that is, show that \(n^{-1} \sum_{i=1}^{n} e_{i}=0\)
2.15. Correlation between Residuals and Explanatory Variables. Consider
 a generic sequence of pairs of numbers \(( x_{1} , \, y_{1} ) , \, \dots, \, ( x_{n} , \, y_{n} )\) with the cor-
 relation coefficient computed as \(r ( y , \, x )=\left[ ( n \,-\, 1 ) s_{y} s_{x} \right]^{-1} \sum_{i=1}^{n} \left( y_{i} \,-\, \overline{{y}} \right)\)
\(( x_{i} \,-\, \overline{{x}} )\) .
a. Suppose that either \(\overline{{y}}=0 , \overline{{x}}=0\) or both \(\overline{{x}}\) and \(\overline{{y}}=0\) . Then, check
 that \(r ( y , \, x )=0\) implies \(\sum_{i=1}^{n} \, y_{i} \, x_{i} \,=0\) and vice versa.
b. Show that the correlation between the residuals and the explanatory
variables is zero. Do this by using part (a) of Exercise 2.13 to show
that \(\sum_{i=1}^{n} \, x_{i} \, e_{i} \,=\, 0\) and then apply part (a).
c. Show that the correlation between the residuals and fitted values is
zero. Do this by showing that \(\sum_{i \mathop=1}^{n} \, \widehat{y}_{i} \, e_{i} \,=\, 0\) and then apply part (a).
2.16. Correlation and \(t\) -Statistics. Use the following steps to establish a rela-
 tionship between the correlation coefficient and the \(t\) -statistic for the slope.
\(\mathrm{a}\) . Use algebra to check that
\[
R^{2}=1-\frac{n-2} {n-1} \frac{s^{2}} {s_{y}^{2}} .
\]
 b. Use part (a) to establish the following quick computational formula
for \(s\) .
\[
s \:=\: s_{y} \sqrt{( 1 \:-\: r^{2} ) \frac{n \:-\: 1} {n \:-\: 2}} .
\]
c. Use part (b) to show that 
\[
t ( b_{1} )=\sqrt{n-2} \frac{r} {\sqrt{1-r^{2}}} .
\]

\clearpage
58 Basic Linear Regression
Sections 2.6-- 2.7
 2.17. Effects of an Unusual Point. You are analyzing a data set of size \(n \,=\, 1 0 0\)
 You have just performed a regression analysis using one predictor variable
 and notice that the residual for the 10th observation is unusually large.
a. Suppose that, in fact, it turns out that \(e_{1 0}=8 s\) . What percentage of the
error sum of squares, Error \(S \, S\) , is due to the 10th observation?
b. Suppose that \(e_{1 0}=4 s\) What percentage of the error sum of squares.
Error \(S \, S\) is due to the 10th observation?
c. Suppose that you reduce the dataset to size \(n=2 0\) After running the
regression, it turns out that we still have \(1 0=4 s\) . What percentage of
the error sum of squares, Error \(\mathit{S S}\) , is due to the 10th observation?
2.18. Consider a dataset consisting of 20 observations with the following sum-
mary statistics: \(\overline{{x}}=0 , \; \overline{{y}}=9 , \; s_{x}=1\) , and \(s_{y}=1 0\) You run a regression
 using using one variable and determine that \(s=7 .\) Determine the standard
error of a prediction at \(x_{*}=1\)
@ EMPIRICAL 2.19. Summary Statistics Can Hide Important Relationships. The data in
Filename is Table 2.9 is due to Anscombe (1973). The purpose of this exercise is to
"AnscombesData" demonstrate how plotting data can reveal important information that is not
evident in numerical summary statistics.
 a. Compute the averages and standard deviations of each column of data.
Check that the averages and standard deviations of each of the \(x\)
columns are the same, within two decimal places, and similarly for
each of the \(y\) columns.
b. Run four regressions, (i) \(y_{1}\) on \(x_{1}\) (ii) \(y_{2}\) on \(x_{1}\) , (iii) \(y_{3}\) on \(x_{1}\) , and
\(( \mathrm{i v} ) ~ y_{4} ~ \mathrm{o n} ~ x_{2}\) . Verify, for each of the four regressions fits, that \(b_{0} \approx3 . 0\)
\(b_{1} \approx0 . 5 , \, s \approx1 . 2 3 7\) , and \(R^{2} \approx0 . 6 7 7 ,\) within two decimal places.
c. Produce scatter plots for each of the four regression models that you
fit in part (b).
d. Discuss the fact that the fitted regression models produced in part
(b) imply that the four datasets are similar, though the four scatter plots
 produced in part (c) yield a dramatically different story.
Table 2.9
Anscombe's (1973) Obs
Data num \(x_{1}\) \(y_{1}\) \(y_{2}\) \(y_{3}\) \(x_{2}\) \(y_{4}\)
1 10 8.04 9.14 7.46 8 6.58
2 8 6.95 8.14 6.77 8 5.76
3 13 7.58 8.74 12.74 8 7.71
4 9 8.81 8.77 7.11 8 8.84
5 11 8.33 9.26 7.81 8 8.47
6 14 9.96 8.10 8.84 8 7.04
7 6 7.24 6.13 6.08 8 5.25
8 4 4.26 3.10 5.39 8 5.56
9 12 10.84 9.13 8.15 8 7.91
10 7 4.82 7.26 6.42 8 6.89
11 5 5.68 4.74 5.73 19 12.50
\clearpage
2.10 Exercises 59
2.20. Nursing Home Utilization. This exercise considers nursing home data@ EMPiRicAL
provided by the Wisconsin Department of Health and Family Services Filename is
(DHFS) and described in Exercises 1.2 and 2.10. "WiscNursingHome"
You decide to examine the relationship between total patient years
(LOGTPY) and the number of beds (LOGNUMBED), both in logarith-
mic units, using cost-report year 2001 data.
a. Summary statistics. Create basic summary statistics for each variable
 Summarize the relationship through a correlation statistic and a scatter
 plot.
b. Fit the basic linear model. Cite the basic summary statistics,
include the coefficient of determination, the regression coefficient for
LOGNUMBED, and the corresponding \(t\) -statistic.
c. Hypothesis testing. Test the following hypotheses at the \(5\) level
of significance using a \(t\) -statistic. Also compute the corresponding
\(p\) -value.
\(\mathrm{c ( i )}\) . Test \(H_{0} : \beta_{1}=0\) versus \(H_{a} : \beta_{1} \neq0 .\)
c(ii). Test \(H_{0} : \beta_{1}=1\) versus \(H_{a} : \beta_{1} \neq1 .\)
\(\mathrm{c}\) (iii). Test \(H_{0} : \beta_{1}=1\) versus \(H_{a} : \beta_{1} > 1 .\)
\(\mathrm{c ( i v )}\) . Test \(H_{0} : \beta_{1}=1\) versus \(H_{a} : \beta_{1} < 1\) .
d. You are interested in the effect that a marginal change in
LOGNUMBED has on the expected value of LOGTPY.
\(\mathrm{d ( i )}\) .Suppose that there is a marginal change in LOGNUMBED of 2.
Provide a point estimate of the expected change in LOGTPY.
d(ii). Provide a 95\% confidence interval corresponding to the point
estimate in part \(\mathrm{d ( i )}\) .
d(ii). Provide a 99\% confidence interval corresponding to the point
estimate in part \(\mathrm{d ( i )}\) .
e. At a specified number of beds estimate \(x_{*}=1 0 0\) do these things:
\(\mathrm{e ( i )}\) .Find the predicted value of LOGTPY.
e(ii). Obtain the standard error of the prediction.
e(iii). Obtain a 95\% prediction interval for your prediction.
\(\mathrm{e ( i v )}\) .: Convert the point prediction in part \(\mathrm{e ( i )}\) and the prediction inter-
val obtained in part e(iii) into total person years (through expo-
nentiation).
\(\mathrm{e ( v )}\) .Obtain a prediction interval as in part \(\mathrm{e ( i v )}\) , corresponding to a
 90\% level (in lieu of 95\%). @ EMPIRICAL
2.21. Initial Public Offerings. As a financial analyst, you want to convince a Filename is
client of the merits of investing in firms that have just entered a stock "IPO"
exchange, as an initial public offering (IPO). Thus, you gather data on 116
firms that priced during the six-month time frame of January 1 through
 June 1, 1998. By looking at this recent historical data, you are able to
 compute RETURN, the firm's one-year return (as a percentage).
You are also interested in looking at financial characteristics of the firm
 that may help you understand (and predict) the return. You initially examine
 REVENUE, the firm's 1997 revenues in millions of dollars. Unfortunately,
\clearpage
60 Basic Linear Regression
Table 2.10 Initial
Public Offering Standard
 Summary Statistics Mean Median Deviation Minimum Maximum
RETURN 0.106 -0.130 0.824 0.938 4.333
REV 134.487 39.971 261.881 0.099 1455.761
LnREV 3.686 3.688 1.698 -2.316 7.283
PRICEIPO 13.195 13.000 4.694 4.000 29.000
this variable was not available for six firms. Thus, the statistics here are
 for the 110 firms that have both REVENUES and RETURNS. In addition,
Table 2.10 provides information on the (natural) logarithmic revenues,
denoted as LnREV, and the initial price of the stock, denoted as PRICEIPO.
a. You hypothesize that larger firms, as measured by revenues, are more
stable and thus should enjoy greater returns. You have determined that
the correlation between RETURN and REVENUE is \(- 0 . 0 1 7 5 .\)
\(\mathrm{a ( i )}\) .Calculate the least squares fit using REVENUE to predict
RETURN. Determine \(b_{0}\) and \(b_{1}\) .
a(ii). For Hyperion Telecommunications, REVENUE is 95.55 (mil-
 lions of dollars). Calculate the fitted RETURN using the regres-
 sion fit in part \(\mathrm{a ( i )}\) .
b. Table 2.11 summarizes a regression using logarithmic revenues and
returns.
b(i). Suppose instead that you use LnREVs to predict RETURN.
Calculate the fitted RETURN under this regression model. Is
this equal your answer in part a(ii)?
 b(ii). Do logarithmic revenues significantly affect returns? To this end,
 provide a formal test of hypothesis. State your null and alterna-
tive hypotheses, decision-making criterion, and your decision-
 making rule. Use a 10\% level of significance.
b(iii). You conjecture that, other things equal, that firms with greater
revenues will be more stable and thus enjoy a larger initial return.
Thus, you wish to consider the null hypothesis of no relation
 between LnREV and RETURN versus the alternative hypothesis
 that there is a positive relation between LnREV and RETURN.
To this end, provide a formal test of hypothesis. State your
null and alternative hypotheses, decision-making criterion, and
 decision-making rule. Use a 10\% level of significance.
 c.Determine the correlation between LnREV and RETURN. Be sure to
state whether the correlation is positive, negative, or zero.
 d. You are considering investing in a firm that has LnRE \(\mathrm{V}=2\) (so revenues
are \(e^{2}=7 . 3 8 9\) millions of dollars).
d(i). Using the fitted regression model, determine the least squares
point prediction.
d(ii). Determine the 95\% prediction interval corresponding to your
 prediction in part \(\mathrm{d ( i )}\) .
\clearpage
2.10 Exercises 61
Table 2.11
Standard Regression Results
 Variable Coefficient Error \(t\) -Statistic from a Model Fit with
Logarithmic
INTERCEPT 0.438 0.186 2.35 Revenues
LnREV \(- 0 . 0 9 0\) 0.046 \(- 1 . 9 7\)
Notes: \(s \,=\, 0 . 8 1 3 6\) and \(R^{2}=0 . 0 3 4 5 2\)
e. The \(R^{2}\) from the fitted regression model is a disappointing 3.5\%. Part of 
 the difficulty is due to observation number 59, the Inktomi Corporation.
Inktomi sales are 12th smallest of the data set, with \(\mathrm{L n R E V}=1 . 7 6\)
(so revenues are \(e^{1 . 7 6}=5 . 7 9\) millions of dollars), yet it has the highest
first-year return, with \(\mathrm{R E T U R N}=4 3 3 . 3 3 .\)
\(\mathrm{e ( i )}\) . Calculate the residual for this observation.
e(ii). What proportion of the unexplained variability (error sum of
 squares) does this observation account for?
\(\mathrm{\Theta( i i i )}\) Define the idea of a high leverage observation.
\(\mathrm{e ( i v )}\) .: Would this observation be considered a high leverage observa-
tion? Justify your answer. @ EMPIRICAL
2.22. National Life Expectancies. We continue the analysis begun in Exercise Filename is
1.7 by examining the relation between \(\boldsymbol{y}=\boldsymbol{L} \boldsymbol{I} \boldsymbol{F} \boldsymbol{E} \boldsymbol{X} \boldsymbol{P}\) and \(x \,=\, F E R T I L I T Y\) ."UNLifeExpectancy"
 shown in Figure 2.12. Fit a linear regression model of LIFEEXP using the
 explanatory variable \(x \,=\, F E R T I L I T Y .\)
: Figure 2.12 Plot of 
FERTILITY versus
\(\S\) T 8 \(\omicron_{\infty}^{\circ}\) C LIFEEXP.
2 0 Q 8
1
88
: 0 0
00:0
o 0
3 : \(\omicron^{\circ}\) .0.0
1
1 2 3 \(4 ~\) 5 6 7
FERTILITY
 a. The United States has a FERTILITY rate of 2.0. Determine the fitted
life expectancy.
b. The island nation Dominica did not report a FERTILITY rate and thus
was not included in the regression. Suppose that its FERTILITY rate
is 2.0. Provide a 95\% prediction interval for the life expectancy in
Dominica.
c. China has a FERTILITY rate of 1.7 and a life expectancy of 72.5
 Determine the residual under the model. How many multiples of \(s\) is
 this residual from zero?
\clearpage
62 Basic Linear Regression
d.Suppose that your prior hypothesis is that the FERTILITY slope is --6.0
 and you wish to test the null hypothesis that the slope has increased
(i.e., the slope is greater than \(- 6 . 0\) . Test this hypothesis at the \(5 7_{o}\)
 level of significance. Also compute an approximate \(p\) -value.
 2.11 Technical Supplement -- Elements of Matrix Algebra
Examples are an excellent tool for introducing technical topics such as regression.
However, this chapter has also used algebra as well as basic probability and
statistics to give you further insights into regression analysis. Going forward,
we will be studying multivariate relationships. With many things happening
concurrently in several dimensions, algebra is no longer useful for providing
insights. Instead, we will need matrix algebra. This supplement provides a brief
 introduction to matrix algebra to allow you to study the linear regression chapters
 of this text. It re-introduces basic linear regression to give you a feel for things that
will be coming up in subsequent chapters when we extend basic linear regression
 to the multivariate case. Appendix A3 defines additional matrix concepts.
 2.11.1 Basic Definitions
A matrix is a rectangular array of numbers arranged in rows and columns (the
 plural of matrix is matrices). For example, consider the income and age of three
people:
\[
\mathbf{A}=\begin{array} {c c} {\quad C o l \ 1} & {C o l \ 2} \\ {R o w \ 1} \\ {R o w \ 2 \left( \begin{matrix} {6 , 0 0 0} & {2 3} \\ {1 3 , 0 0 0} & {4 7} \\ {1 1 , 0 0 0} & {3 5} \\ \end{matrix} \right)} \\ \end{array} .
\]
Here, column \(1\) represents income, and column \(2\) represents age. Each row
corresponds to an individual. For example, the first individual is 23 years old
with an income of \$6,000.
The number of rows and columns is called the dimension of the matrix. For
example, the dimension of the matrix \(\mathbf{A}\) above is \(3 \times2\) (read \(3 \mathrm{\iota\iota\ b y \iota\iota\iota\iota}\) . This
stands for three rows and two columns. If we were to represent the income and
 age of 100 people, then the dimension of the matrix would be \(1 0 0 \times2\) .
It is convenient to represent a matrix using the notation
\[
\mathbf{A}=\left( \begin{matrix} {a_{1 1}} & {a_{1 2}} \\ {a_{2 1}} & {a_{2 2}} \\ {a_{3 1}} & {a_{3 1}} \\ \end{matrix} \right) .
\]
Here, \(a_{i j}\) is the symbol for the number in the ith row and \(j\) th column of \({\bf A}\) . In
general, we work with matrices of the form
\[
\mathbf{A}=\left( \begin{matrix} {a_{1 1}} & {a_{1 2}} & {\cdots} & {a_{1 k}} \\ {\vdots} & {\vdots} & {\ddots} & {\vdots} \\ {a_{n 1}} & {a_{n 2}} & {\cdots} & {a_{n k}} \\ \end{matrix} \right) .
\]
In this case, the matrix A has dimension \(n \times k\) .
\clearpage
2.11 Technical Supplement - Elements of Matrix Algebra 63
A vector is a special matrix. A row vector is a matrix containing only 1 row
( \(n=1\) ). A column vector is a matrix containing only one column ( \(k=1\) . For
example,
\[
\mathrm{c o l u m n ~ v e c t o r} \to\left( \begin{matrix} {2} \\ {3} \\ {4} \\ {5} \\ {6} \\ \end{matrix} \right) \quad\mathrm{r o w ~ v e c t o r} \to\left( \begin{matrix} {2} & {3} & {4} & {5} & {6} \\ \end{matrix} \right) .
\]
Notice that the row vector takes much less room on a printed page than the
corresponding column vector. A basic operation that relates these two quantities
 is the transpose. The transpose of a matrix A is defined by interchanging the rows
 and columns and is denoted by \({\mathbf{A}^{\prime}} \ ( \mathrm{o r} \ {\mathbf{A}^{T}} )\) ). For example,
\[
\begin{array} {c c} {\mathbf{A}=\left( \begin{matrix} {6 0 0 0} & {2 3} \\ {1 3 {,} 0 0 0} & {4 7} \\ {1 1 {,} 0 0 0} & {3 5} \\ \end{matrix} \right)} & {\mathbf{A}^{\prime}=\left( \begin{matrix} {6 0 0 0} & {1 3 {,} 0 0 0} & {1 1 {,} 0 0 0} \\ {2 3} & {4 7} & {3 5} \\ \end{matrix} \right) .} \\ \end{array}
\]
Thus, if \({\bf A}\) has dimension \(n \times k\) then \(\mathbf{A^{\prime}}\) has dimension \(k \times n\) .
2.11.2 Some Special Matrices
(i) A square matrix is a matrix where the number of rows equals the number
of columns; that is, \(n=k\) .
(ii) The diagonal numbers of a square matrix are the numbers of a matrix
 where the row number equals the column number, for example, \(a_{1 1} , \, a_{2 2}\) ,
 and so on. A diagonal matrix is a square matrix where all nondiagonal
numbers are equal to \(0\) . For example,
\[
\mathbf{A}=\left( \begin{matrix} {-1} & {0} & {0} \\ {0} & {2} & {0} \\ {0} & {0} & {3} \\ \end{matrix} \right) .
\]
(iii) An identity matrix is a diagonal matrix where all the diagonal numbers
 are equal to 1. This special matrix is often denoted by \(\mathbf{I}\)
(iv) A symmetric matrix is a square matrix \(\mathbf{A}\) such that the matrix remains
unchanged if we interchange the roles of the rows and columns. More
formally, a matrix \({\bf A}\) is symmetric if \(\mathbf{A}=\mathbf{A^{\prime}}\) . For example,
\[
{\bf A}=\left( \begin{matrix} {1} & {2} & {3} \\ {2} & {4} & {5} \\ {3} & {5} & {1 0} \\ \end{matrix} \right)={\bf A}^{\prime} .
\]
Note that a diagonal matrix is a symmetric matrix.
 2.11.3 Basic Operations
 Scalar Multiplication
Let \({\bf A}\) be a \(n \times k\) matrix and let \(c\) be a real number. That is, a real number is a
\(1 \times1\) matrix and is also called a scalar. Multiplying a scalar \(c\) by a matrix \(\mathbf{A}\) is
\clearpage
 64 Basic Linear Regression
denoted by \(c {\bf A}\) and defined by
\[
c \mathbf{A}=\left( \begin{matrix} {c a_{1 1}} & {c a_{1 2}} & {\cdots} & {c a_{1 k}} \\ {\vdots} & {\vdots} & {\ddots} & {\vdots} \\ {c a_{n 1}} & {c a_{n 2}} & {\cdots} & {c a_{n k}} \\ \end{matrix} \right) .
\]
For example, suppose that \(c=1 0\) and
\[
\mathbf{A}=\left( \! \! \begin{array} {c c} {1} & {2} \\ {6} & {8} \\ \end{array} \! \! \right) \quad\mathrm{t h e n} \quad\mathbf{B}=c \mathbf{A}=\left( \! \! \begin{array} {c c} {1 0} & {2 0} \\ {6 0} & {8 0} \\ \end{array} \! \! \right) .
\]
 Note that \(c \mathbf{A}=\mathbf{A} c\) .
 Addition and Subtraction of Matrices
Let \(\mathbf{A}\) and \({\bf B}\) be matrices with dimensions \(n \times k\) . Use \(a_{i j}\) and \(b_{i j}\) to denote the
numbers in the ith row and \(j\) th column of \(\mathbf{A}\) and \({\bf B}\) , respectively. Then, the
matrix \(\mathbf{C}=\mathbf{A}+\mathbf{B}\) is defined to be the matrix with \(( a_{i j}+b_{i j} )\) in the \(i\) th row and
\(j\) th column. Similarly, the matrix \(\mathbf{C}=\mathbf{A}-\mathbf{B}\) is defined to be the matrix with
\(( a_{i j} \,-\, b_{i j} )\) in the \(i\) th row and \(j\) th column. Symbolically, we write this as the
following:
\[
\mathrm{I f} \ \ \ \mathbf{A}=\left( a_{i j} \right)_{i j} \ \ \ \mathrm{a n d} \ \ \ \mathbf{B}=\left( b_{i j} \right)_{i j} \, , \, \mathrm{t h e n}
\]
\[
\mathbf{C}=\mathbf{A}+\mathbf{B}=\left( a_{i j} \,+\, b_{i j} \right)_{i j} \ \ \mathrm{~ a n d ~} \ \ \mathbf{C}=\mathbf{A}-\mathbf{B}=\left( a_{i j} \,-\, b_{i j} \right)_{i j} .
\]
For example, consider
\[
\mathbf{A}=\left( \! \! \begin{array} {c c} {2} & {5} \\ {4} & {1} \\ \end{array} \! \! \right) \ \ \ \ \mathbf{B}=\left( \! \! \begin{array} {c c} {4} & {6} \\ {8} & {1} \\ \end{array} \! \! \right) .
\]
Then
\[
\mathbf{A}+\mathbf{B}=\left( \! \! \begin{array} {c c} {6} & {1 1} \\ {1 2} & {2} \\ \end{array} \! \! \right) \ \ \ \ \mathbf{A}-\mathbf{B}=\left( \! \! \begin{array} {c c} {-2} & {-1} \\ {-4} & {0} \\ \end{array} \! \! \right) .
\]
Basic Linear Regression Example of Addition and Subtraction. Now, recall
 that the basic linear regression model can be written as \(n\) equations:
\[
y_{1}=\beta_{0}+\beta_{1} x_{1}+\varepsilon_{1}
\]
:
\[
y_{n}=\beta_{0}+\beta_{1} x_{n}+\varepsilon_{n} .
\]
We can define
\[
\mathbf{y}=\left( \begin{matrix} {y_{1}} \\ {\vdots} \\ {y_{n}} \\ \end{matrix} \right) \quad\boldsymbol{\varepsilon}=\left( \begin{matrix} {\varepsilon_{1}} \\ {\vdots} \\ {\varepsilon_{n}} \\ \end{matrix} \right) \quad\quad\mathrm{a n d} \quad\textrm{E} \mathbf{y}=\left( \begin{matrix} {\beta_{0}+\beta_{1} x_{1}} \\ {\vdots} \\ {\beta_{0}+\beta_{1} x_{n}} \\ \end{matrix} \right) .
\]
With this notation, we can express the \(n\) equations more compactly as \(\mathbf{y} \!=\! \textrm{E} \mathbf{y} \!+\! \boldsymbol{\varepsilon}\) .
\clearpage
 2.11 Technical Supplement - Elements of Matrix Algebra 65
 Matrix Multiplication
In general, if \(\mathbf{A}\) is a matrix of dimension \(n \times c\) and \({\bf B}\) is a matrix of dimension
\(c \times k\) , then \(\mathbf{C}=\mathbf{A} \mathbf{B}\) is a matrix of dimension \(n \times k\) and is defined by
\[
\begin{aligned} {\mathbf{C}=\mathbf{A} \mathbf{B}=} & {{} \left( \sum_{s=1}^{c} a_{i \, s} b_{s j} \right)_{i j} .} \\ \end{aligned}
\]
For example consider the \(2 \times2\) matrices
\[
\mathbf{A}=\left( \begin{matrix} {2} & {5} \\ {4} & {1} \\ \end{matrix} \right) \quad\mathrm{a n d} \quad\mathbf{B}=\left( \begin{matrix} {4} & {6} \\ {8} & {1} \\ \end{matrix} \right) .
\]
The matrix AB has dimension \(2 \times2 .\) To illustrate the calculation, consider the
number in the first row and second column of AB. By the rule presented ear-
lier, with \(i \ =1\) and \(j=2\) the corresponding element of AB is \(\sum_{s=1}^{2} a_{1 s} b_{s 2}=\)
\(a_{1 1} b_{1 2} \,+\, a_{1 2} b_{2 2}=2 ( 6 )+\, 5 ( 1 )=1 7 .\) The other calculations are summarized as
\[
\mathbf{A B}=\left( \begin{matrix} {2 ( 4 )+5 ( 8 )} & {2 ( 6 )+5 ( 1 )} \\ {4 ( 4 )+1 ( 8 )} & {4 ( 6 )+1 ( 1 )} \\ \end{matrix} \right)=\left( \begin{matrix} {4 8} & {1 7} \\ {2 4} & {2 5} \\ \end{matrix} \right) .
\]
As another example, suppose
\[
\mathbf{A}=\left( \begin{matrix} {1} & {2} & {4} \\ {0} & {5} & {8} \\ \end{matrix} \right) \quad\mathbf{B}=\left( \begin{matrix} {3} \\ {5} \\ {2} \\ \end{matrix} \right) .
\]
Because \({\bf A}\) has dimension \(2 \times3\) and \({\bf B}\) has dimension \(3 \times1\) this means that the
 product AB has dimension \(2 \times1\) The calculations are summarized as
\[
\mathbf{A B}=\left( \begin{matrix} {1 ( 3 )+2 ( 5 )+4 ( 2 )} \\ {0 ( 3 )+5 ( 5 )+( 2 )} \\ \end{matrix} \right)=\left( \begin{matrix} {2 1} \\ {4 1} \\ \end{matrix} \right) .
\]
For some additional examples, we have
\[
\left( \begin{matrix} {4} & {2} \\ {5} & {8} \\ \end{matrix} \right) \left( \begin{matrix} {a_{1}} \\ {a_{2}} \\ \end{matrix} \right)=\left( \begin{matrix} {4 a_{1}+2 a_{2}} \\ {5 a_{1}+8 a_{2}} \\ \end{matrix} \right) .
\]
\[
2 \quad3 \quad5 \bigr) \left( \begin{matrix} {2} \\ {3} \\ {5} \\ \end{matrix} \right)=2^{2}+3^{2}+5^{2}=3 8 \quad\ \left( \begin{matrix} {2} \\ {3} \\ {5} \\ \end{matrix} \right) \left( \begin{matrix} {2} & {3} & {5} \\ \end{matrix} \right)=\left( \begin{matrix} {4} & {6} & {1 0} \\ {6} & {9} & {1 5} \\ {1 0} & {1 5} & {2 5} \\ \end{matrix} \right) .
\]
In general, you see that \(\mathbf{A B} \neq\mathbf{B} \mathbf{A}\) in matrix multiplication, unlike multiplication
of scalars (real numbers). Further, we remark that the identity matrix serves
the role of " one"in matrix multiplication, in that \(\mathbf{A I}=\mathbf{A}\) and \(\mathbf{I A}=\mathbf{A}\) for any
matrix \({\bf A}\) providing that the dimensions are compatible to allow matrix multipli-
cation.
\clearpage
66 Basic Linear Regression
 Basic Linear Regression Example of Matrix Multiplication. Define
\[
\mathbf{X}=\left( \! \! \begin{array} {c c} {1} & {x_{1}} \\ {\vdots} & {\vdots} \\ {1} & {x_{n}} \\ \end{array} \! \! \right) \ \ \mathrm{a n d} \ \ \boldsymbol{\beta}=\left( \! \! \begin{array} {c} {\beta_{0}} \\ {\beta_{1}} \\ \end{array} \! \! \right) , \mathrm{t o ~ g e t ~} \mathbf{X} \boldsymbol{\beta}=\left( \! \! \begin{array} {c} {\beta_{0}+\beta_{1} x_{1}} \\ {\vdots} \\ {\beta_{0}+\beta_{1} x_{n}} \\ \end{array} \! \! \right)=\mathrm{E} \ \mathbf{y} .
\]
Thus, this yields the familiar matrix expression of the regression model,
\(\mathbf{y}=\mathbf{X} \boldsymbol{\beta}+\boldsymbol{\varepsilon}\) . Other useful quantities include
\[
\mathbf{y}^{\prime} \mathbf{y}=\left( \! \! \begin{array} {c c c} {y_{1}} & {\cdots} & {y_{n}} \\ \end{array} \! \! \right) \left( \! \! \begin{array} {c} {y_{1}} \\ {\vdots} \\ {y_{n}} \\ \end{array} \! \! \right)=y_{1}^{2}+\cdots+y_{n}^{2}=\sum_{i=1}^{n} y_{i}^{2} ,
\]
\[
\mathbf{X}^{\prime} \mathbf{y}=\left( \! \! \begin{array} {c c c} {1} & {\cdots} & {1} \\ {x_{1}} & {\cdots} & {x_{n}} \\ \end{array} \! \! \right) \left( \! \! \begin{array} {c} {y_{1}} \\ {\vdots} \\ {y_{n}} \\ \end{array} \! \! \right)=\left( \! \! \begin{array} {c} {\sum_{i=1}^{n} y_{i}} \\ {\sum_{i=1}^{n} x_{i} \, y_{i}} \\ \end{array} \! \! \right) ,
\]
and
\[
\mathbf{X}^{\prime} \mathbf{X}=\left( \begin{matrix} {1} & {\cdots} & {1} \\ {x_{1}} & {\cdots} & {x_{n}} \\ \end{matrix} \right) \left( \begin{matrix} {1} & {x_{1}} \\ {\vdots} & {\vdots} \\ {1} & {x_{n}} \\ \end{matrix} \right)=\left( \begin{matrix} {n} & {\sum_{i=1}^{n} x_{i}} \\ {\sum_{i=1}^{n} x_{i}} & {\sum_{i=1}^{n} x_{i}^{2}} \\ \end{matrix} \right) .
\]
Note that \(\mathbf{X}^{\prime} \mathbf{X}\) is a symmetric matrix.
Matrix Inverses
In matrix algebra, there is no concept of " division." Instead, we extend the concept
of " reciprocals"of real numbers. To begin, suppose that \(\mathbf{A}\) is a square matrix of
dimension \(k \times k\) and let \(\mathbf{I}\) be the \(k \times k\) identity matrix. If there exists a \(k \times k\)
matrix \({\bf B}\) such that \(\mathbf{A} \mathbf{B}=\mathbf{I}=\mathbf{B} \mathbf{A}\) , then \({\bf B}\) is called the inverse of \(\mathbf{A}\) and is written
\[
{\bf B}={\bf A}^{-1} .
\]
Now, not all square matrices have inverses. Further, even when an inverse exists,
 it may not be easy to compute by hand. One exception to this rule are diagonal
matrices. Suppose that \({\bf A}\) is diagonal matrix of the form
\[
\mathbf{A}=\left( \! \! \begin{array} {c c c} {a_{1 1}} & {\cdots} & {0} \\ {\vdots} & {\ddots} & {\vdots} \\ {0} & {\cdots} & {a_{k k}} \\ \end{array} \! \! \right) . \; \; \; \mathrm{T h e n} \; \; \; \mathbf{A}^{-1} \!=\! \left( \! \! \begin{array} {c c c} {{\frac{1} {a_{1 1}}}} & {\cdots} & {0} \\ {\vdots} & {\ddots} & {\vdots} \\ {0} & {\cdots} & {{\frac{1} {a_{k k}}}} \\ \end{array} \! \! \right) .
\]
For example,
\[
\begin{array} {c c c c} {\left( \begin{matrix} {2} & {0} \\ {0} & {-1 9} \\ \end{matrix} \right)} & {\left( \begin{matrix} {\frac{1} {2}} & {0} \\ {0} & {-\frac{1} {1 9}} \\ \end{matrix} \right)} & {=} & {\left( \begin{matrix} {1} & {0} \\ {0} & {1} \\ \end{matrix} \right)} \\ {\bf{A}} & {\bf{A}^{-1}} & {=} & {\bf{I}} \\ \end{array} .
\]

\clearpage
2.11 Technical Supplement - Elements of Matrix Algebra 67
In the case of a matrix of dimension \(2 \times2\) , the inversion procedure can be
 accomplished by hand easily even when the matrix is not diagonal. In the \(2 \times2\)
case, we suppose that if
\[
\mathbf{A}=\left( \! \! \begin{array} {c c} {a} & {b} \\ {c} & {d} \\ \end{array} \! \! \right) , \quad\mathrm{t h e n} \quad\mathbf{A}^{-1} {=} \frac{1} {a d \,-\, b c} \left( \! \! \begin{array} {c c} {d} & {-b} \\ {-c} & {a} \\ \end{array} \! \! \right) .
\]
Thus, for example, if
\[
\mathbf{A}=\left( \begin{matrix} {2} & {2} \\ {3} & {4} \\ \end{matrix} \right) \ \ \ \mathrm{t h e n} \ \ \ \mathbf{A}^{-1} {=} \frac{1} {2 ( 4 )-2 ( 3 )} \left( \begin{matrix} {4} & {-2} \\ {-3} & {2} \\ \end{matrix} \right)=\left( \begin{matrix} {2} & {-1} \\ {-3 / 2} & {1} \\ \end{matrix} \right) .
\]
As a check, we have
\[
{\bf A} {\bf A}^{-1}=\left( {\begin{matrix} {2} & {2} \\ {3} & {4} \\ \end{matrix}} \right) \left( {\begin{matrix} {2} & {-1} \\ {-3 / 2} & {1} \\ \end{matrix}} \right)=\left( {2 ( 2 )-2 ( 3 / 2 )} & {2 (-1 )+2 ( 1 )} \\ {3 ( 2 )-4 ( 3 / 2 )} & {3 (-1 )+4 ( 1 )} \\ \end{matrix} \right)=\left( {\begin{matrix} {1} & {0} \\ {0} & {1} \\ \end{matrix}} \right)={\bf I} .
\]
Basic Linear Regression Example of Matrix Inverses. With
\[
\mathbf{X}^{\prime} \mathbf{X}=\! \left( \begin{matrix} {n} & {\sum_{i=1}^{n} x_{i}} \\ {\sum_{i=1}^{n} x_{i}} & {\sum_{i=1}^{n} x_{i}^{2}} \\ \end{matrix} \right) ,
\]
we have
\[
\left( \mathbf{X}^{*} \mathbf{X} \right)^{-1} \!=\! \! \frac{1} {n \, \sum_{i=1}^{n} x_{i}^{2}-\left( \sum_{i=1}^{n} x_{i} \right)^{2}} \left( \begin{matrix} {\sum_{i=1}^{n} x_{i}^{2}} & {-\sum_{i=1}^{n} x_{i}} \\ {-\sum_{i=1}^{n} x_{i}} & {n} \\ \end{matrix} \right) .
\]
To simplify this expression, recall that \(\overline{{x}}=n^{-1} \sum_{\substack{i=1}}^{n} x_{i}\) . Thus,
\[
{\bf( X^{\prime} {\bf X} )}^{-1} \!=\! \frac{1} {\sum_{i=1}^{n} x_{i}^{2}-n \overline{{x}}^{2}} \left( \begin{matrix} {n^{-1} \sum_{i=1}^{n} x_{i}^{2}} & {-\overline{{x}}} \\ {-\overline{{x}}} & {1} \\ \end{matrix} \right) . \tag{2.9}
\]
 Section 3.1 will discuss the relation \(\mathbf{b}=\left( \mathbf{X}^{\prime} \mathbf{X} \right)^{-1} \mathbf{X}^{\prime} \mathbf{y}\) To illustrate the calcu-
lation, we have
\[
\mathbf{b}=\left( \mathbf{X}^{\prime} \mathbf{X} \right)^{-1} \mathbf{X}^{\prime} \mathbf{y}=\! \! \frac{1} {\sum_{i=1}^{n} x_{i}^{2}-n \overline{{x}}^{2}} \left( \! \! \begin{array} {c c} {n^{-1} \sum_{i=1}^{n} x_{i}^{2}} & {-\overline{{x}}} \\ {-\overline{{x}}} & {1} \\ \end{array} \! \! \right) \left( \! \! \begin{array} {c} {\sum_{i=1}^{n} y_{i}} \\ {\sum_{i=1}^{n} x_{i} y_{i}} \\ \end{array} \! \! \right)
\]
\[
= \frac{1} {\sum_{i=1}^{n} x_{i}^{2}-n \overline{x}^{2}} \left( \begin{aligned} {\sum_{i=1}^{n} \left( \overline{y} x_{i}^{2}-\overline{x} x_{i} \, y_{i} \right)} \\ {\sum_{i=1}^{n} x_{i} \, y_{i}-n \overline{x} \, y} \\ \end{aligned} \right)=\left( \begin{matrix} {b_{0}} \\ {b_{1}} \\ \end{matrix} \right) .
\]

\clearpage
68 Basic Linear Regression
From this expression, we may see
\[
b_{1}=\frac{\sum_{i=1}^{n} x_{i} \, y_{i} \,-n \overline{{x}} \, \overline{{y}}} {\sum_{i=1}^{n} x_{i}^{2}-n \overline{{x}}^{2}} ,
\]
and
\[
b_{0}=\frac{\overline{{y}} \sum_{i=1}^{n} x_{i}^{2}-\overline{{x}} \sum_{i=1}^{n} x_{i} \, y_{i}} {\sum_{i=1}^{n} x_{i}^{2}-n \overline{{x}}^{2}}=\frac{\overline{{y}} \left( \sum_{i=1}^{n} x_{i}^{2}-n \overline{{x}}^{2} \right)-\overline{{x}} \left( \sum_{i=1}^{n} x_{i} \, y_{i}-n \overline{{x}} \, y \right)} {\sum_{i=1}^{n} x_{i}^{2}-n \overline{{x}}^{2}}=\overline{{y}}-b_{1} \overline{{x}} .
\]
These are the usual expressions for the slope \(b_{1}\) (Exercise 2.8) and intercept \(b_{0}\) .
 2.11.4 Random Matrices
Expectations. Consider a matrix of random variables
\[
{\bf U}=\left( \begin{matrix} {u_{1 1}} & {u_{1 2}} & {\cdots} & {u_{1 c}} \\ {u_{2 1}} & {u_{2 2}} & {\cdots} & {u_{2 c}} \\ {\vdots} & {\vdots} & {\ddots} & {\vdots} \\ {u_{n 1}} & {u_{n 2}} & {\cdots} & {u_{n c}} \\ \end{matrix} \right) .
\]
When we write the expectation of a matrix, this is shorthand for the matrix
 of expectations. Specifically, suppose that the joint probability function of \(u_{1 1}\)
\(u_{1 2} , \, \ldots, \, u_{1 c} , \, \ldots, \, u_{n 1} , \, \ldots, \, u_{n c}\) is available to define the expectation operator.
Then we define
\[
\textrm{E} \mathbf{U}=\left( \begin{matrix} {\mathrm{E} u_{1 1}} & {\mathrm{E} u_{1 2}} & {\cdots} & {\mathrm{E} u_{1 c}} \\ {\mathrm{E} u_{2 1}} & {\mathrm{E} u_{2 2}} & {\cdots} & {\mathrm{E} u_{2 c}} \\ {\vdots} & {\vdots} & {\ddots} & {\vdots} \\ {\mathrm{E} u_{n 1}} & {\mathrm{E} u_{n 2}} & {\cdots} & {\mathrm{E} u_{n c}} \\ \end{matrix} \right) .
\]
As an important special case, consider the joint probability function for the
 random variables \(y_{1} , \dots, \, y_{n}\) and the corresponding expectations operator. Then
\[
\textrm{E} \mathbf{y}=\mathrm{E} \left( \begin{matrix} {y_{1}} \\ {\vdots} \\ {y_{n}} \\ \end{matrix} \right)=\left( \begin{matrix} {\textrm{E} y_{1}} \\ {\vdots} \\ {\textrm{E} y_{n}} \\ \end{matrix} \right) .
\]
\(\mathrm{B y}\) the linearity of expectations, for a nonrandom matrix \(\mathbf{A}\) and vector \({\bf B}\) , we have
\(\textrm{E} ( \mathbf{A} \textbf{y}+\mathbf{B} )=\mathbf{A} \textrm{E} \textbf{y}+\mathbf{B}\) .
Variances. We can also work with second moments of random vectors. The
 variance of a vector of random variables is called the variance-covariance matrix.
It is defined by
\[
\mathrm{V a r ~} \mathbf{y}=\mathrm{E} ( ( \mathbf{y}-\mathrm{E} \, \mathbf{y} ) ( \mathbf{y}-\mathrm{E} \, \mathbf{y} )^{\prime} ) . \tag{2.10}
\]

\clearpage
2.11 Technical Supplement - Elements of Matrix Algebra 69
That is, we can express
\[
\begin{aligned} {\mathrm{V a r ~} \mathbf{y}} & {{}=\mathrm{E} \left( \left( \begin{matrix} {y_{1}-\mathrm{E} \, y_{1}} \\ {\vdots} \\ {y_{n}-\mathrm{E} \, y_{n}} \\ \end{matrix} \right) \bigl( y_{1}-\mathrm{E} \, y_{1}} & {\cdots} & {y_{n}-\mathrm{E} \, y_{n} \, \bigr) \right)} \\ {} & {{}=\left( \begin{matrix} {\mathrm{V a r ~} y_{1}} & {\mathrm{C o v} ( y_{1} , \, y_{2} )} & {\cdots} & {\mathrm{C o v} ( y_{1} , \, y_{n} )} \\ {\mathrm{C o v} ( y_{2} , \, y_{1} )} & {\mathrm{V a r ~} y_{2}} & {\cdots} & {\mathrm{C o v} ( y_{2} , \, y_{n} )} \\ {\vdots} & {\vdots} & {\ddots} & {\vdots} \\ {\mathrm{C o v} ( y_{n} , \, y_{1} )} & {\mathrm{C o v} ( y_{n} , \, y_{2} )} & {\cdots} & {\mathrm{V a r ~} y_{n}} \\ \end{matrix} \right) ,} \\ \end{aligned}
\]
because \(\mathrm{E} ( ( y_{i} \,-\mathrm{E} \, y_{i} ) ( y_{j} \,-\mathrm{E} \, y_{j} ) )=\mathrm{C o v} ( y_{i} , \, y_{j} )\) for \(i \neq j\) and \(\mathrm{C o v} ( y_{i} , \, y_{i} )=\)
Var \(y_{i}\) .
In the case that \(y_{1} , \dots, \, y_{n}\) are mutually uncorrelated, we have that \(\mathrm{C o v} ( y_{i}\)
\(y_{j} )=0\) for \(i \neq j\) and thus
\[
\mathrm{V a r ~} \mathbf{y}=\! \left( \! \! \begin{array} {c c c c} {\mathrm{V a r ~} y_{1}} & {0} & {\cdots} & {0} \\ {0} & {\mathrm{V a r ~} y_{2}} & {\cdots} & {0} \\ {\vdots} & {\vdots} & {\ddots} & {\vdots} \\ {0} & {0} & {\cdots} & {\mathrm{V a r ~} y_{n}} \\ \end{array} \! \! \right) .
\]
Further, if the variances are identical so that Var \(y_{i}=\sigma^{2}\) then we can write
Var \(\mathbf{y}=\sigma^{2} \mathbf{I}\) where \(\mathbf{I}\) is the \(n \times n\) identity matrix. For example, if \(y_{1} , \dots, \, y_{n}\) are
i.i.d., then Var \(\mathbf{y}=\sigma^{2} \mathbf{I} .\)
From equation (2.10), it can be shown that
\[
\mathrm{V a r} \left( {\bf A y+B} \right)=\mathrm{V a r} \left( {\bf A y} \right)={\bf A} \left( {\mathrm{V a r} \; {\bf y}} \right) {\bf A^{\prime}} . \tag{2.11}
\]
For example, if \({' {\bf A}=( a_{1} , \, a_{2} , \, \ldots, \, a_{n} )}={\bf a}^{\prime} \, \mathrm{a n d} \, {\bf B=} \, {\bf0} ,\) then equation (2.11) reduces
to
\[
\begin{aligned} {\mathrm{V a r} \left( \sum_{i=1}^{n} a_{i} \, y_{i} \right)} & {{}=\mathrm{V a r} \left( \mathbf{a}^{\prime} \, \mathbf{y} \right)=\mathbf{a}^{\prime} \left( \mathrm{V a r} \, \mathbf{y} \right) \mathbf{a}=\left( a_{1} , \, a_{2} , \, \dots, \, a_{n} \right) \left( \mathrm{V a r} \, \mathbf{y} \right) \left( \begin{matrix} {a_{1}} \\ {\vdots} \\ {a_{n}} \\ \end{matrix} \right)} \\ {} & {{}=\sum_{i=1}^{n} a_{i}^{2} \mathrm{V a r} \, \, y_{i} \, \,+\, \, 2 \sum_{i=2}^{n} \sum_{j=1}^{i-1} a_{i} a_{j} \mathrm{C o v} ( y_{i} , \, y_{j} ) .} \\ \end{aligned}
\]
 Defnition -- Multivariate Normal Distribution. A vector of random variables
\(\mathbf{y}=( y_{1} , \, \ldots, \, y_{n} )^{\prime}\) is said to be multivariate normal if all linear combinations of
 the form \(\sum_{i=1}^{n} a_{i} \, y_{i}\) are normally distributed. In this case, we write \({\bf y} \sim N ( {\boldsymbol\mu} , \, \Sigma) ,\)
where \(\boldsymbol{\mu}=\operatorname{E} {\bf y}\) is the expected value of \({\bf y}\) and \(\Sigma=\mathrm{V a r} \, {\bf y}\) is the variance-covariance
matrix of \({\bf y}\) . From the definition, we have that \(\mathbf{y} \sim N ( \boldsymbol{\mu} , \, \Sigma)\) implies that \(\mathbf{a}^{\prime} \mathbf{y} \sim\)
mean \(\mu\sum_{i=1}^{n} a_{i}\) and variance \(\sigma^{2} \sum_{i=1}^{n} a_{i}^{2}\) \(\sum_{i=1}^{n} a_{i} \, y_{i}\) is distributed normally with
\(N ( \mathbf{a}^{\prime} \boldsymbol{\mu} , \, \mathbf{a}^{\prime} \boldsymbol{\Sigma} \mathbf{a} )\) Thus, if \(y_{i}\) are i.i.d., then
\clearpage

\end{document}
